\documentclass[11pt]{article}
\usepackage[margin=1in]{geometry}
\usepackage{enumitem}
% =============================================================================
% Preambles/header.tex — shared LaTeX/Beamer preamble for this project
%
% Use from a Beamer lecture by prepending:
%     \input{header}
% and compiling with TEXINPUTS=../Preambles:$TEXINPUTS (the /compile-latex
% skill does this automatically).
%
% PALETTE CONTRACT. Color names defined here MUST match the names in
% Quarto/theme-template.scss so Beamer and Quarto renderings use the same
% palette. The scripts/check-palette-sync.sh script enforces this.
%
% To customize: edit the HEX values below AND the matching $variable in
% Quarto/theme-template.scss, then run scripts/check-palette-sync.sh.
% =============================================================================

% -----------------------------------------------------------------------------
% Engine detection + encoding
%
% Our /compile-latex skill uses XeLaTeX, where inputenc is unnecessary and
% can produce encoding warnings. Only load inputenc under pdfTeX.
% -----------------------------------------------------------------------------
\usepackage{iftex}
\ifPDFTeX
  \usepackage[utf8]{inputenc}
\fi

\usepackage{amsmath, amssymb, amsthm}
\usepackage{booktabs}
\usepackage{graphicx}

% -----------------------------------------------------------------------------
% PALETTE — mirrors Quarto/theme-template.scss variable names.
% Load xcolor and define colors BEFORE anything that references them
% (notably hyperref, hypersetup, and the Beamer color assignments).
% -----------------------------------------------------------------------------
\usepackage{xcolor}

\definecolor{primary-blue}{HTML}{012169}      % headings, accents
\definecolor{primary-gold}{HTML}{B9975B}      % emphasis, borders
\definecolor{highlight-yellow}{HTML}{F2A900}  % markers, alerts
\definecolor{light-bg}{HTML}{E8EDF5}          % title / section backgrounds
\definecolor{jet}{HTML}{1A1A1A}               % body text

% Semantic colors (used in diagrams and annotations)
\definecolor{positive}{HTML}{15803D}          % good / observed / identified
\definecolor{negative}{HTML}{B91C1C}          % bad / problematic / bias
\definecolor{neutral}{HTML}{525252}           % reference / context

% Highlight accents (match .hi-* classes in the SCSS)
\definecolor{hi-slate}{HTML}{314F4F}
\definecolor{hi-green}{HTML}{2E8B57}
\definecolor{hi-red}{HTML}{C41E3A}

% Semantic aliases so skill templates and snippets can write
% \textcolor{accent}{...} without hard-coding "primary-blue".
\colorlet{accent}{primary-blue}
\colorlet{accent2}{primary-gold}

% -----------------------------------------------------------------------------
% Hyperref — loaded AFTER xcolor + palette so link colors resolve.
% -----------------------------------------------------------------------------
\usepackage{hyperref}
\hypersetup{colorlinks=true, linkcolor=primary-blue, urlcolor=primary-blue,
            citecolor=primary-blue, pdfborder={0 0 0}}

% -----------------------------------------------------------------------------
% Course code — set once per deck (after \input{header}, before \begin{document})
% via \coursecode{CS401}. Shown in the footer on every slide so a deck's course
% is identifiable without opening the title page. Empty by default.
% -----------------------------------------------------------------------------
\makeatletter
\newcommand{\coursecode}[1]{\gdef\@coursecodetext{#1}}
\gdef\@coursecodetext{}
\makeatother

% -----------------------------------------------------------------------------
% Beamer theme assignments (only apply under Beamer).
%
% We detect Beamer via \@ifclassloaded{beamer}, which is the documented,
% non-internal way. This must be wrapped in \makeatletter/\makeatother
% because \@ifclassloaded contains an @-catcode character.
% -----------------------------------------------------------------------------
\makeatletter
\@ifclassloaded{beamer}{%
  \usefonttheme{professionalfonts}
  \setbeamercolor{structure}{fg=primary-blue}
  \setbeamercolor{frametitle}{fg=primary-blue}
  \setbeamercolor{title}{fg=primary-blue}
  \setbeamercolor{subtitle}{fg=primary-gold}
  \setbeamercolor{author}{fg=primary-blue}
  \setbeamercolor{institute}{fg=neutral}
  \setbeamercolor{date}{fg=primary-gold}
  \setbeamercolor{itemize item}{fg=primary-blue}
  \setbeamercolor{itemize subitem}{fg=primary-gold}
  \setbeamercolor{alerted text}{fg=primary-gold}
  \setbeamercolor{block title}{fg=white, bg=primary-blue}
  \setbeamercolor{block body}{bg=light-bg}
  % Semantic block variants: block=definition/concept (blue), exampleblock=
  % worked example (green/positive), alertblock=key takeaway or caution (gold).
  % Keep to <=2 colored boxes per slide across ALL three variants (INV-7).
  \setbeamercolor{block title example}{fg=white, bg=positive}
  \setbeamercolor{block body example}{bg=light-bg}
  \setbeamercolor{block title alerted}{fg=white, bg=primary-gold}
  \setbeamercolor{block body alerted}{bg=light-bg}

  % Minimal footer: course code (left) + page X / N (right), no navigation clutter
  \setbeamertemplate{navigation symbols}{}
  \setbeamertemplate{footline}{%
    \color{neutral}\footnotesize\hspace{0.7em}\@coursecodetext\hfill
    \insertframenumber/\inserttotalframenumber\hspace{0.7em}\vspace{0.3em}}
}{}
\makeatother

% -----------------------------------------------------------------------------
% TikZ setup — libraries the snippet gallery uses
% -----------------------------------------------------------------------------
\usepackage{tikz}
\usetikzlibrary{arrows.meta, positioning, calc, decorations.pathreplacing,
                fit, shapes.geometric, backgrounds}

% Shared TikZ styles — snippets and hand-written diagrams can pick these up
% via `\begin{tikzpicture}[every node/.style={...}]` or by naming specific
% styles (e.g., `\node[dag-node] (X) at (0,0) {X};`).
\tikzset{
  % Node styles for causal DAGs and similar
  dag-node/.style = {
    draw=primary-blue, fill=light-bg, rounded corners=2pt,
    minimum width=1.2cm, minimum height=0.9cm, align=center, font=\small
  },
  decision-node/.style = {
    draw=primary-gold, fill=white, diamond, aspect=1.8,
    minimum width=1.8cm, minimum height=1.2cm, align=center, font=\small,
    inner sep=1pt
  },
  % Edge styles — observed vs counterfactual
  observed-edge/.style = {-{Latex[length=2.5mm]}, thick, draw=jet},
  counterfactual-edge/.style = {-{Latex[length=2.5mm]}, thick, draw=neutral, dashed},
  confound-edge/.style = {-{Latex[length=2.5mm]}, thick, draw=negative, dashed},
  % Data-point styles for plots
  observed-dot/.style = {fill=primary-blue, draw=primary-blue, circle, inner sep=1.2pt},
  counterfactual-dot/.style = {fill=white, draw=neutral, circle, inner sep=1.2pt}
}

% -----------------------------------------------------------------------------
% Convenience macros
% -----------------------------------------------------------------------------
\newcommand{\muted}[1]{\textcolor{neutral}{#1}}
\newcommand{\key}[1]{\textcolor{primary-gold}{\textbf{#1}}}
\newcommand{\good}[1]{\textcolor{positive}{#1}}
\newcommand{\bad}[1]{\textcolor{negative}{#1}}

% Standout frame macro: full-bleed dark background for section transitions.
% Beamer-only (uses \begin{frame}/\setbeamercolor/\usebeamerfont) -- do not
% call from an article-class file (e.g. Notes/<CODE>/*-notes.tex). \muted,
% \key, \good, \bad, and \coursecode above are plain xcolor/gdef and are
% safe to use from either class; verified when Notes/article-class support
% was added.
% Expand: \transitionslide{Your transition title}
\newcommand{\transitionslide}[1]{%
  \begingroup
  \setbeamercolor{background canvas}{bg=primary-blue}
  \begin{frame}[plain]
    \centering
    \vfill
    {\usebeamerfont{title}\color{white}\Large #1\par}
    \vfill
  \end{frame}
  \endgroup
}

% Section divider frame (Beamer-only), an alternative to \transitionslide with a
% darker two-line style: near-black (jet) background, a small white label line
% above a larger gold title, and a thin gold rule along the bottom edge.
% Expand: \sectiondivider{Section 2}{Pointers}
\newcommand{\sectiondivider}[2]{%
  \begingroup
  \setbeamercolor{background canvas}{bg=jet}
  \begin{frame}[plain]
    \vspace*{3.0cm}
    {\color{white}\ttfamily\large #1\par}
    \vspace{0.5cm}
    {\color{primary-gold}\LARGE #2\par}
    \begin{tikzpicture}[remember picture, overlay]
      \draw[primary-gold, line width=2.5pt]
        ($(current page.south west)+(0,0.1)$) -- ($(current page.south east)+(0,0.1)$);
    \end{tikzpicture}
  \end{frame}
  \endgroup
}

% -----------------------------------------------------------------------------
% Channel watermark — "Concepts That Click" (YouTube).
%
% Article-class files (CompetitiveExam Q/A sets, Notes, Assignments) get a
% light diagonal draft watermark on every page plus a centered footer naming
% the channel. Beamer decks get a subtle TikZ background watermark instead
% (the footline already carries the course code). Centralized here so every
% MCQ PDF is branded without per-file edits.
% -----------------------------------------------------------------------------
\makeatletter
\@ifclassloaded{article}{%
  \usepackage{draftwatermark}
  \SetWatermarkText{Concepts That Click}
  \SetWatermarkScale{1}
  \SetWatermarkFontSize{2cm}
  \SetWatermarkLightness{0.86}
  \SetWatermarkAngle{45}
  \usepackage{fancyhdr}
  \pagestyle{fancy}
  \fancyhf{}
  \fancyfoot[C]{\footnotesize\textcolor{neutral}{Concepts That Click~|~YouTube: \href{https://www.youtube.com/@concepts.thatclick}{@concepts.thatclick}}}
  \renewcommand{\headrulewidth}{0pt}
  \renewcommand{\footrulewidth}{0pt}
  \fancypagestyle{plain}{%
    \fancyhf{}%
    \fancyfoot[C]{\footnotesize\textcolor{neutral}{Concepts That Click~|~YouTube: \href{https://www.youtube.com/@concepts.thatclick}{@concepts.thatclick}}}%
    \renewcommand{\headrulewidth}{0pt}%
    \renewcommand{\footrulewidth}{0pt}%
  }
}{}
\@ifclassloaded{beamer}{%
  \setbeamertemplate{background}{%
    \begin{tikzpicture}[remember picture, overlay]
      \node[rotate=45, opacity=0.055, align=center]
        at (current page.center)
        {\Huge\textcolor{neutral}{Concepts That Click}};
    \end{tikzpicture}%
  }
}{}
\makeatother 
\coursecode{JKSSB}
\title{Lesson 08 Practice: Answers and Rationales}
\author{JKSSB Computer Awareness}
\date{\today}
\begin{document}
\maketitle
\noindent\textbf{Provenance:} All 20 items are \textit{[Original]}
JKSSB-pattern questions for this lesson. No item is a real PYQ, and the
answer key is the author's verification of the lesson's concepts.
\begin{enumerate}[label=\textbf{Q\arabic*.}, leftmargin=*]
\item \textbf{C.} RAM is volatile, so it loses its contents when power is
  removed. ROM, a hard disk, and a flash drive are non-volatile and retain
  their contents.
\item \textbf{A.} RAM stands for Random Access Memory. B, C, and D are
  invented expansions.
\item \textbf{A.} ROM stands for Read-Only Memory. B, C, and D are invented
  expansions.
\item \textbf{C.} RAM is volatile and does \emph{not} retain its contents
  after power-off, so C is the statement that is NOT correct. A, B, and D
  are all true statements.
\item \textbf{B.} RAM is the volatile main memory holding running programs
  and their data. ROM is non-volatile; L1 cache is smaller working memory;
  the SSD is secondary storage.
\item \textbf{A.} DRAM is a volatile form of RAM. ROM, PROM, and EPROM
  belong to the non-volatile ROM family.
\item \textbf{D.} SRAM is faster than DRAM and needs no refresh, which is
  why it is used for cache. DRAM needs periodic refresh; ROM and a hard disk
  are non-volatile and slower than SRAM.
\item \textbf{D.} DRAM is denser and cheaper than SRAM but needs periodic
  refresh. A describes SRAM; B is the use of SRAM; C is false, since DRAM is
  volatile.
\item \textbf{B.} SRAM is used for cache memory because it is fast and needs
  no refresh. A is DRAM's role; C and D describe other memory classes.
\item \textbf{A.} DRAM serves as main memory: dense and cheap, at the cost of
  periodic refresh. B is SRAM's role; C describes ROM; D is a register, not
  DRAM.
\item \textbf{B.} EEPROM is Electrically Erasable Programmable Read-Only
  Memory. A should be Dynamic, C should be Erasable (UV), and D should be
  Static, so each of the others is a wrong expansion.
\item \textbf{D.} PROM is programmable once and cannot be erased. EPROM and
  EEPROM can be erased; DRAM is a volatile RAM, not a ROM type.
\item \textbf{B.} EPROM is erased by ultraviolet light. PROM cannot be
  erased, EEPROM is erased electrically, and mask ROM is fixed at
  manufacture.
\item \textbf{C.} EEPROM is erased and rewritten electrically, in-circuit.
  PROM cannot be erased, EPROM needs UV light plus removal, and SRAM is
  volatile RAM.
\item \textbf{B.} Flash is an EEPROM-derived non-volatile memory erased in
  blocks. A wrongly makes it volatile DRAM; C describes a fixed ROM type; D
  is a register, not flash.
\item \textbf{A.} Statements I and II are correct: RAM is volatile and ROM
  is non-volatile, and ROM is normally read-only though specially
  programmable. III is false because cache is faster than main memory, not
  slower. B and C each include false III; D includes all three.
\item \textbf{A.} L1 is the smallest, fastest cache and sits closest to the
  core. L2 and L3 are larger and slower; main memory is slower still.
\item \textbf{D.} L3 is the largest cache level and is typically shared
  across cores. L1 and L2 are smaller and faster; ROM is not cache.
\item \textbf{C.} A miss means the data is not found in the cache, so it is
  fetched from main memory. A describes a hit; B and D are unrelated.
\item \textbf{D.} Hit ratio $=$ hits $\div$ total accesses $= 90/100 = 0.90$.
  A is the miss fraction, B would mean every access hit, and C is
  arithmetically wrong.

\item \textbf{A.} RAM is volatile and readable/writable. B, C, and D wrongly make it permanent or tie it to peripherals.
\item \textbf{A.} Firmware lives in non-volatile ROM/Flash. RAM (B) loses contents; C and D are not firmware stores.
\item \textbf{A.} SRAM's speed suits cache. Tape (B), floppy coating (C), and paper (D) are not cache technologies.
\item \textbf{A, B, C.} RAM, cache, and registers lose contents without power. Masked ROM (D) is non-volatile.
\item \textbf{A, B, C.} DRAM refresh, SRAM speed, and SRAM-based cache are correct. D is false: ROM keeps contents without power.
\end{enumerate}
\def\JKSSBLessonID{8}\def\JKSSBAnswers{1}\newcommand{\JKSSBItem}[5]{%
  \ifnum\JKSSBAnswers=1
    \item \textbf{#4.} #5
  \else
    \item \textit{#1} #2
      \begin{enumerate}[label=\Alph*.]
        #3
      \end{enumerate}
  \fi
}

\begin{enumerate}[label=\textbf{Q\arabic*.}, leftmargin=*]
\ifnum\JKSSBLessonID=50
  \setcounter{enumi}{26}
\else\ifnum\JKSSBLessonID=59
  \setcounter{enumi}{27}
\else
  \setcounter{enumi}{25}
\fi\fi
\ifnum\JKSSBLessonID=1
\JKSSBItem{[Original]}{In the IPO cycle, which stage converts processed data into information for a user?}
  {\item Input
   \item Processing
   \item Output
   \item Storage}
  {C}{Output presents the processed result. Input supplies data, processing transforms it, and storage retains it.}
\JKSSBItem{[Original]}{Which property describes a computer's ability to repeat a task for long periods without tiring?}
  {\item Diligence
   \item Versatility
   \item Accuracy
   \item Automation}
  {A}{Diligence means sustained work without fatigue. Versatility is range of uses, accuracy concerns correct results, and automation is operation without continuous manual control.}
\JKSSBItem{[Original]}{A list of unprocessed daily sales entries is best described as:}
  {\item Information
   \item Data
   \item Knowledge
   \item Output}
  {B}{Raw facts are data. Information is data organized or processed into a meaningful form; knowledge and output are not synonyms for raw entries.}
\JKSSBItem{[Original]}{Which statement is a limitation of a computer?}
  {\item It can process data rapidly.
   \item It can store large amounts of data.
   \item It does not independently decide its goals without instructions.
   \item It can perform repetitive work consistently.}
  {C}{A computer follows instructions and does not set its own goals. The other options describe speed, storage, and diligence.}
\JKSSBItem{[Original, MSQ drill]}{\textbf{(MSQ: select ALL correct options.)} Which statements about the IPO cycle are correct?}
  {\item Input provides data or instructions.
   \item Processing applies operations to data.
   \item Output presents a result.
   \item Storage is the same stage as input.}
  {A, B, C}{Input, processing, and output have distinct roles. Storage retains data or results; it is not the same as input.}
\fi
\ifnum\JKSSBLessonID=2
\JKSSBItem{[Original]}{A computer that represents values using continuously varying physical quantities is:}
  {\item Digital
   \item Analog
   \item Hybrid
   \item Microcomputer}
  {B}{Analog computers represent continuous quantities. Digital computers use discrete values; hybrid systems combine analog and digital processing; microcomputer classifies size/use.}
\JKSSBItem{[Original]}{A system combining analog measurement with digital control is a:}
  {\item Hybrid computer
   \item Mainframe
   \item Minicomputer
   \item General-purpose digital computer}
  {A}{A hybrid computer combines analog and digital characteristics. The other options identify size/use categories or a digital-only class.}
\JKSSBItem{[Original]}{Which class is designed to support very large numbers of users and high-volume transaction processing?}
  {\item Mainframe
   \item Tablet computer
   \item Embedded controller
   \item Analog computer}
  {A}{Mainframes serve many concurrent users and transaction workloads. Tablets are personal devices, embedded controllers serve dedicated functions, and analog describes data representation.}
\JKSSBItem{[Original]}{Which classification is based primarily on the kind of data represented rather than physical size?}
  {\item Analog, digital, hybrid
   \item Micro, mini, mainframe
   \item Laptop, desktop, tablet
   \item Client, server, workstation}
  {A}{Analog/digital/hybrid describes representation. The remaining sets classify by size/form factor or role.}
\JKSSBItem{[Original, MSQ drill]}{\textbf{(MSQ: select ALL correct options.)} Which pairings are valid?}
  {\item Analog computer --- continuous quantities
   \item Digital computer --- discrete values
   \item Hybrid computer --- combination of analog and digital features
   \item Mainframe --- a type of analog representation}
  {A, B, C}{The first three match data-representation categories. Mainframe describes a computer class by scale and workload, not representation.}
\fi
\ifnum\JKSSBLessonID=3
\JKSSBItem{[PYQ-adapted: JKSSB Website Operator 2026 Q72, provisional key; SRC-04/SRC-07]}{What is the smallest unit of data?}
  {\item Bit
   \item Byte
   \item Nibble
   \item Word}
  {A}{A bit is a single binary digit. A nibble is four bits, a byte is eight bits, and a word is processor-dependent. The cited answer key is provisional.}
\JKSSBItem{[Original]}{How many bits are in one byte?}
  {\item 4
   \item 8
   \item 16
   \item 32}
  {B}{One byte contains eight bits. Four is a nibble; 16 and 32 bits are larger groups.}
\JKSSBItem{[Original]}{How many nibbles make one byte?}
  {\item One
   \item Two
   \item Four
   \item Eight}
  {B}{A nibble is four bits and a byte is eight bits, so one byte contains two nibbles.}
\JKSSBItem{[Original]}{Using the binary convention, how many bytes are in 1 KiB?}
  {\item 1000
   \item 1024
   \item 100
   \item 2048}
  {B}{1 KiB is 1024 bytes. 1000 is the decimal kilobyte convention; 100 and 2048 are not the stated unit sizes.}
\JKSSBItem{[Original, MSQ drill]}{\textbf{(MSQ: select ALL correct options.)} Which statements are correct?}
  {\item A nibble contains four bits.
   \item A byte contains eight bits.
   \item A bit can represent 0 or 1.
   \item A byte contains four nibbles.}
  {A, B, C}{A nibble is four bits and a byte is eight bits, so a byte has two nibbles, not four.}
\fi
\ifnum\JKSSBLessonID=4
\JKSSBItem{[PYQ-adapted: JKSSB Website Operator 2026 Q65, provisional key; SRC-04/SRC-07]}{Which unit directs the sequence of operations and controls data flow inside the CPU?}
  {\item Arithmetic Logic Unit
   \item Control Unit
   \item Main memory
   \item Input unit}
  {B}{The Control Unit sequences operations and directs data flow. The ALU performs arithmetic/logic; memory stores data; input supplies data. The cited key is provisional.}
\JKSSBItem{[PYQ-adapted: JKSSB Junior Assistant 2026 Q61, final key; SRC-03/SRC-06]}{Which CPU component performs arithmetic and logical operations on binary data?}
  {\item Control Unit
   \item Arithmetic Logic Unit
   \item Program Counter
   \item Input unit}
  {B}{The ALU performs arithmetic and logical operations. The Control Unit directs execution, the Program Counter holds an instruction address, and the input unit supplies data.}
\JKSSBItem{[Original]}{Which bus carries the address of a memory location selected by the CPU?}
  {\item Address bus
   \item Data bus
   \item Control bus
   \item Expansion bus}
  {A}{The address bus carries location addresses. The data bus carries data; the control bus carries timing/control signals; expansion buses connect peripherals.}
\JKSSBItem{[Original]}{Which sequence best describes a basic computer system's information flow?}
  {\item Output, processing, input
   \item Input, processing, output
   \item Storage, output, input
   \item Processing, output, input}
  {B}{Input supplies data, processing transforms it, and output presents results. The other sequences reverse that basic order.}
\JKSSBItem{[Original, MSQ drill]}{\textbf{(MSQ: select ALL correct options.)} Which are functional units of a general-purpose computer?}
  {\item Input unit
   \item Memory unit
   \item Central Processing Unit
   \item Output unit}
  {A, B, C, D}{A general-purpose computer has input, memory, processing (CPU), and output functions.}
\fi
\ifnum\JKSSBLessonID=5
\JKSSBItem{[PYQ-adapted: SSC CGL 2025 Tier-II Paper-I, 19-Jan-2026 Q6; SRC-18]}{Which register holds the instruction currently being decoded?}
  {\item Memory Address Register
   \item Program Counter
   \item Instruction Register
   \item Accumulator}
  {C}{The Instruction Register holds the fetched instruction during decode/execute. MAR holds an address, PC holds the next instruction address, and the accumulator stores intermediate results.}
\JKSSBItem{[PYQ-adapted: JKSSB Junior Assistant 2026 Q64, final key; SRC-03/SRC-06]}{Which register normally holds the address of the next instruction to be fetched?}
  {\item Instruction Register
   \item Program Counter
   \item Accumulator
   \item Memory Data Register}
  {B}{The Program Counter holds the next instruction address. The IR holds the instruction itself, the accumulator holds results, and MDR holds data transferred to/from memory.}
\JKSSBItem{[Original]}{Which register temporarily holds the address of a memory location being accessed?}
  {\item MAR
   \item IR
   \item Accumulator
   \item Status register}
  {A}{The Memory Address Register holds a memory address. IR holds an instruction, the accumulator stores results, and a status register stores condition flags.}
\JKSSBItem{[Original]}{The clock rate of a processor is commonly expressed in:}
  {\item Bytes per second
   \item Hertz
   \item Pixels per inch
   \item Revolutions per minute}
  {B}{Clock rate is cycles per second, measured in hertz. The other units measure data rate, image density, and rotational speed.}
\JKSSBItem{[Original, MSQ drill]}{\textbf{(MSQ: select ALL correct options.)} Match each register with its common role.}
  {\item PC --- address of the next instruction
   \item IR --- current instruction being decoded
   \item MAR --- memory address for a transfer
   \item Accumulator --- intermediate arithmetic/logic result}
  {A, B, C, D}{All four register-role pairings are correct.}
\fi
\ifnum\JKSSBLessonID=6
\JKSSBItem{[Original]}{Which processor is generally optimized for a wide range of sequential general-purpose tasks?}
  {\item CPU
   \item GPU
   \item Network interface card
   \item Sound card}
  {A}{A CPU is a general-purpose processor. GPUs specialize in parallel graphics/data workloads; a NIC and sound card are peripheral controllers.}
\JKSSBItem{[Original]}{A graphics workload applies the same operation to many pixels in parallel. Which component is commonly suited to this work?}
  {\item GPU
   \item Hard disk
   \item Keyboard controller
   \item ROM}
  {A}{A GPU has many parallel processing elements suited to graphics operations. Storage, input control, and firmware memory do not perform this computation.}
\JKSSBItem{[Original]}{A microcontroller is most likely to combine a CPU with:}
  {\item Only an external display
   \item Memory and input/output peripherals on one chip
   \item A separate supercomputer
   \item A printer interface only}
  {B}{A microcontroller integrates processing, memory, and peripherals for embedded control. The other choices do not describe that integrated design.}
\JKSSBItem{[Original]}{Which is a typical example of an embedded-system use for a microcontroller?}
  {\item Dedicated washing-machine control
   \item A large-scale weather simulation
   \item A multi-user banking database server
   \item A general-purpose desktop operating system}
  {A}{A washing-machine controller performs a dedicated embedded task. The other workloads typically use general-purpose or high-performance computer systems.}
\JKSSBItem{[Original, MSQ drill]}{\textbf{(MSQ: select ALL correct options.)} Which statements are generally correct?}
  {\item CPUs are designed for broad general-purpose processing.
   \item GPUs can execute many similar operations in parallel.
   \item Microcontrollers commonly serve embedded control tasks.
   \item A GPU is another name for a hard disk.}
  {A, B, C}{The first three describe processor roles. A GPU processes data; a hard disk stores it.}
\fi
\ifnum\JKSSBLessonID=7
\JKSSBItem{[PYQ-adapted: JKSSB Junior Assistant 2026 Q63, final key; SRC-03/SRC-06]}{Which memory is typically faster than main memory and holds copies of recently used data?}
  {\item Cache
   \item HDD
   \item Optical disc
   \item Magnetic tape}
  {A}{Cache sits between CPU and RAM to reduce access time. HDDs, optical discs, and tape are secondary storage and are slower.}
\JKSSBItem{[Original]}{In the usual speed hierarchy, which is fastest?}
  {\item Registers
   \item RAM
   \item SSD
   \item Magnetic tape}
  {A}{Registers are closest to the CPU and fastest. RAM, SSD, and tape are progressively slower storage levels.}
\JKSSBItem{[Original]}{As capacity increases and cost per bit generally falls, memory in a computer is usually:}
  {\item Faster at every level
   \item Slower at each lower level of the hierarchy
   \item Volatile at every level
   \item Removed from the CPU}
  {B}{Larger, cheaper storage is generally slower, which motivates the hierarchy. Not all levels are volatile, and memory is not removed from the CPU.}
\JKSSBItem{[Original]}{Which pair is correctly ordered from faster to slower?}
  {\item RAM, registers, cache
   \item Registers, cache, RAM
   \item SSD, RAM, registers
   \item Tape, cache, registers}
  {B}{The usual order is registers, cache, then RAM. The other options reverse one or more levels.}
\JKSSBItem{[Original, MSQ drill]}{\textbf{(MSQ: select ALL correct options.)} Which statements describe the memory hierarchy?}
  {\item Registers are very small and fast.
   \item Cache can reduce average access time.
   \item Secondary storage is generally non-volatile.
   \item Magnetic tape is faster than CPU registers.}
  {A, B, C}{Registers are small/fast, cache reduces access time, and secondary storage retains data without power. Tape is much slower than registers.}
\fi
\ifnum\JKSSBLessonID=8
\JKSSBItem{[PYQ-adapted: JKSSB Sub-Inspector 2022, QCRIAS item 391; SRC-17/SRC-22]}{What normally happens to data held only in RAM when power is switched off?}
  {\item It is retained permanently.
   \item It is lost.
   \item It is copied automatically to ROM.
   \item It is stored in the CPU cache.}
  {B}{RAM is volatile, so its contents are lost when power is removed. It is not automatically copied to ROM or cache.}
\JKSSBItem{[PYQ-adapted: JKSSB FAA 2022, QCRIAS item 551; SRC-17/SRC-21]}{Which memory type is traditionally associated with storing startup firmware such as BIOS?}
  {\item RAM
   \item ROM
   \item CPU cache
   \item Register file}
  {B}{Firmware must be available at startup and stored in non-volatile memory, traditionally ROM. The other listed memories are not the standard firmware store.}
\JKSSBItem{[Original]}{Which memory type is volatile?}
  {\item RAM
   \item ROM
   \item SSD
   \item DVD}
  {A}{RAM loses its contents without power. ROM, SSDs, and optical discs are non-volatile.}
\JKSSBItem{[Original]}{Which type of RAM is commonly built from flip-flop circuits and does not need periodic refreshing?}
  {\item DRAM
   \item SRAM
   \item ROM
   \item Flash memory}
  {B}{SRAM uses flip-flop cells and does not require periodic refresh. DRAM uses capacitors and must be refreshed; ROM and flash are non-volatile.}
\JKSSBItem{[Original, MSQ drill]}{\textbf{(MSQ: select ALL correct options.)} Which statements are correct?}
  {\item RAM is normally volatile.
   \item ROM is non-volatile.
   \item DRAM requires periodic refresh.
   \item SRAM is slower than DRAM in the memory hierarchy.}
  {A, B, C}{RAM is normally volatile, ROM retains data, and DRAM needs refresh. SRAM is generally faster than DRAM, not slower.}
\fi
\ifnum\JKSSBLessonID=9
\JKSSBItem{[PYQ-adapted: JKSSB Junior Assistant 2026 Q62, final key; SRC-03/SRC-06]}{Virtual memory uses part of which resource to extend the apparent capacity of main memory?}
  {\item Secondary storage
   \item CPU register file
   \item Printer buffer
   \item ROM firmware area}
  {A}{Virtual memory uses storage space, usually on an SSD/HDD, as backing for memory pages. Registers, printer buffers, and firmware ROM do not provide this backing store.}
\JKSSBItem{[Original]}{A page currently used by a process is absent from RAM but present in its backing store. The operating system typically responds with a:}
  {\item Page fault
   \item Cache hit
   \item Power failure
   \item Print spool}
  {A}{A page fault signals that the referenced page is not resident in RAM and must be fetched. A cache hit, power failure, and print spool are different events.}
\JKSSBItem{[Original]}{Virtual memory allows a process to address a space that can be:}
  {\item Larger than the available physical RAM
   \item Stored only in CPU registers
   \item Volatile only when the disk is disconnected
   \item Equal to the size of the cache}
  {A}{Virtual addressing can provide a larger apparent address space than physical RAM. Registers and cache do not define the process's full virtual address space.}
\JKSSBItem{[Original]}{Which statement best distinguishes virtual memory from cache?}
  {\item Virtual memory extends the apparent main-memory space; cache speeds up frequent access.
   \item Both are permanent backups.
   \item Cache replaces the operating system.
   \item Virtual memory is a CPU register.}
  {A}{Virtual memory uses storage as backing for memory pages, while cache keeps copies of data for speed. Neither is a permanent backup, OS replacement, or register.}
\JKSSBItem{[Original, MSQ drill]}{\textbf{(MSQ: select ALL correct options.)} Which statements about virtual memory are correct?}
  {\item It is managed with operating-system support.
   \item A page fault can occur when a needed page is not in RAM.
   \item Backing storage may be used for non-resident pages.
   \item It makes secondary storage as fast as registers.}
  {A, B, C}{The OS manages virtual memory, faults locate missing pages, and backing storage can hold non-resident pages. Storage remains much slower than registers.}
\fi
\ifnum\JKSSBLessonID=10
\JKSSBItem{[PYQ-adapted: JKSSB Website Operator 2026 Q67, provisional key; SRC-04/SRC-07]}{Which listed secondary storage device generally has the highest transfer rate?}
  {\item Magnetic tape
   \item HDD
   \item SSD
   \item CD-ROM}
  {C}{An SSD has no moving read/write head and generally transfers data faster than HDDs, tape, or optical discs. The cited key is provisional.}
\JKSSBItem{[PYQ-adapted: JKSSB FAA 2022, QCRIAS item 550; SRC-21]}{Hard-disk rotational speed is commonly specified in:}
  {\item RPM
   \item Pixels per inch
   \item Hertz of the CPU clock
   \item Bits per byte}
  {A}{RPM means revolutions per minute. The alternatives measure display density, CPU frequency, and data grouping.}
\JKSSBItem{[Original]}{Which medium is optical storage?}
  {\item DVD
   \item SSD
   \item HDD
   \item Magnetic tape}
  {A}{A DVD stores data optically. SSDs use flash, HDDs use magnetic platters, and tape is magnetic sequential storage.}
\JKSSBItem{[Original]}{Which storage medium is most commonly associated with sequential access and archival backup?}
  {\item Magnetic tape
   \item CPU cache
   \item RAM
   \item Register}
  {A}{Magnetic tape is a removable sequential-access medium commonly used for backup/archive. The others are volatile CPU/main memory.}
\JKSSBItem{[Original, MSQ drill]}{\textbf{(MSQ: select ALL correct options.)} Which are non-volatile secondary storage media?}
  {\item SSD
   \item HDD
   \item DVD
   \item CPU register}
  {A, B, C}{SSDs, HDDs, and DVDs retain data without power. CPU registers are volatile internal storage.}
\fi
\ifnum\JKSSBLessonID=11
\JKSSBItem{[Original]}{Which backup type copies all selected data each time it runs?}
  {\item Full backup
   \item Incremental backup
   \item Differential backup
   \item Transaction log only}
  {A}{A full backup copies the entire selected set. Incremental copies changes since the latest backup; differential copies changes since the latest full backup.}
\JKSSBItem{[Original]}{After a full backup, an incremental backup normally copies files changed since:}
  {\item The previous backup of any kind
   \item The previous full backup only
   \item The computer was purchased
   \item The last restore test}
  {A}{Incremental backups record changes since the most recent backup. Differential backups use the last full backup as their reference point.}
\JKSSBItem{[Original]}{Which backup usually grows cumulatively until the next full backup?}
  {\item Differential
   \item Incremental
   \item Mirror deletion
   \item Temporary cache}
  {A}{Each differential backup includes changes since the last full backup, so the set can grow. Incrementals reset their comparison after each backup.}
\JKSSBItem{[Original]}{The 3-2-1 backup rule recommends keeping:}
  {\item Three copies, on two media types, with one copy off-site
   \item Three files, on one device, with two deleted
   \item Two copies only, both online
   \item One copy in RAM and two in cache}
  {A}{The rule calls for three copies, two different media, and one off-site copy. The other choices do not provide this redundancy.}
\JKSSBItem{[Original, MSQ drill]}{\textbf{(MSQ: select ALL correct options.)} Which practices improve backup resilience?}
  {\item Keep an off-site copy.
   \item Test restoration from backups.
   \item Keep at least one copy isolated from routine access.
   \item Treat RAID as a complete substitute for backup.}
  {A, B, C}{Off-site, tested, and isolated copies improve recoverability. RAID improves availability but does not protect against deletion, corruption, or every disaster.}
\fi
\ifnum\JKSSBLessonID=12
\JKSSBItem{[PYQ-adapted: JKSSB Junior Assistant 2026 Q68, final key; SRC-03/SRC-06]}{Which device is used to enter hand-drawn strokes directly into a computer?}
  {\item Digitizer tablet
   \item Plotter
   \item Projector
   \item Speaker}
  {A}{A digitizer tablet is an input device. A plotter, projector, and speaker present output.}
\JKSSBItem{[Original]}{Which device captures still images or video as computer input?}
  {\item Webcam
   \item Plotter
   \item Projector
   \item Printer}
  {A}{A webcam captures visual input. Plotters, projectors, and printers are output devices.}
\JKSSBItem{[Original]}{A user moves a pointer by rolling a ball mounted on a stationary pointing device. Which device is this?}
  {\item Trackball
   \item Joystick
   \item Scanner
   \item Light pen}
  {A}{A trackball moves the pointer by rotating a ball while the device remains stationary. The other devices use a stick, capture images, or detect a screen position.}
\JKSSBItem{[Original]}{Which input device is most suitable for recording spoken instructions?}
  {\item Microphone
   \item Monitor
   \item Plotter
   \item Speaker}
  {A}{A microphone converts sound to input. A monitor, plotter, and speaker deliver output.}
\JKSSBItem{[Original, MSQ drill]}{\textbf{(MSQ: select ALL correct options.)} Which are input devices?}
  {\item Scanner
   \item Keyboard
   \item Microphone
   \item Plotter}
  {A, B, C}{Scanner, keyboard, and microphone provide input. A plotter produces output.}
\fi
\ifnum\JKSSBLessonID=13
\JKSSBItem{[Original]}{Which technology converts printed characters into editable text?}
  {\item OCR
   \item OMR
   \item MICR
   \item MIDI}
  {A}{OCR recognizes printed characters as text. OMR detects marked responses, MICR reads magnetic-ink characters, and MIDI describes music data.}
\JKSSBItem{[PYQ-adapted: JKSSB Website Operator 2026 Q70, provisional key; SRC-04/SRC-07]}{Which technology is used to read characters printed in magnetic ink on bank cheques?}
  {\item OMR
   \item OCR
   \item MICR
   \item QR recognition}
  {C}{MICR reads magnetic-ink characters. OMR detects marks, OCR recognizes ordinary printed text, and QR recognition reads two-dimensional codes. The cited key is provisional.}
\JKSSBItem{[PYQ-adapted: JKSSB FAA 2022, QCRIAS item 553; SRC-21]}{A scanned page must be converted into editable text. Which technology is the best fit?}
  {\item OCR
   \item OMR
   \item MICR
   \item Barcode printing}
  {A}{OCR recognizes character shapes in an image and produces text. OMR detects marks, MICR reads magnetic ink, and barcode printing creates codes rather than extracting text.}
\JKSSBItem{[Original]}{OMR is particularly useful for:}
  {\item Detecting shaded response bubbles on forms
   \item Reading cheque account numbers in magnetic ink
   \item Converting speech to audio output
   \item Printing a document}
  {A}{OMR detects marked regions, such as answer bubbles. MICR reads cheque characters; the other choices concern audio output and printing.}
\JKSSBItem{[Original, MSQ drill]}{\textbf{(MSQ: select ALL correct options.)} Match each recognition technology to its task.}
  {\item OCR --- printed text recognition
   \item OMR --- mark detection
   \item MICR --- magnetic-ink character reading
   \item OMR --- magnetic-ink character reading}
  {A, B, C}{OCR, OMR, and MICR have distinct roles. Option D assigns MICR's role to OMR.}
\fi
\ifnum\JKSSBLessonID=14
\JKSSBItem{[PYQ-adapted: JKSSB Website Operator 2026 Q69, provisional key; SRC-04/SRC-07]}{Which output device is designed for producing large engineering drawings?}
  {\item Plotter
   \item Joystick
   \item Scanner
   \item Microphone}
  {A}{A plotter produces large-format drawings. The other devices provide input. The cited answer key is provisional.}
\JKSSBItem{[Original]}{A monitor displays a report on screen without making a paper copy. This is:}
  {\item Soft copy output
   \item Hard copy output
   \item Input
   \item Secondary storage}
  {A}{A screen display is soft copy output. Hard copy is printed output; input and storage have different roles.}
\JKSSBItem{[Original]}{Which device provides visual output on a large screen for an audience?}
  {\item Projector
   \item Scanner
   \item Barcode reader
   \item Trackball}
  {A}{A projector displays visual output. The other devices capture input or move a pointer.}
\JKSSBItem{[Original]}{A tactile display primarily presents information through:}
  {\item Raised or refreshable tactile elements
   \item Magnetic ink
   \item Printed barcodes
   \item Audio sampling only}
  {A}{A tactile display conveys information through touch, often with refreshable braille cells. The other options describe unrelated media/encoding.}
\JKSSBItem{[Original, MSQ drill]}{\textbf{(MSQ: select ALL correct options.)} Which are output devices?}
  {\item Monitor
   \item Speaker
   \item Plotter
   \item Scanner}
  {A, B, C}{Monitor, speaker, and plotter deliver output. A scanner captures input.}
\fi
\ifnum\JKSSBLessonID=15
\JKSSBItem{[Original]}{Which printer forms characters by striking an inked ribbon against paper?}
  {\item Dot-matrix printer
   \item Inkjet printer
   \item Laser printer
   \item Thermal printer}
  {A}{A dot-matrix printer is an impact printer that strikes a ribbon. Inkjet sprays droplets, laser uses toner and a drum, and thermal printing uses heat-sensitive media.}
\JKSSBItem{[Original]}{Which printer is generally classified as non-impact?}
  {\item Laser printer
   \item Daisy-wheel printer
   \item Dot-matrix printer
   \item Line printer}
  {A}{A laser printer forms images without striking paper. Daisy-wheel, dot-matrix, and line printers use impact mechanisms.}
\JKSSBItem{[Original]}{Which device produces a hard copy?}
  {\item Printer
   \item Monitor
   \item Speaker
   \item Projector}
  {A}{A printer produces a physical paper copy. The others provide soft-copy visual or audio output.}
\JKSSBItem{[Original]}{A laser printer typically uses which consumable?}
  {\item Toner
   \item Magnetic ink ribbon
   \item Liquid ink cartridge only
   \item Thermal paper roll only}
  {A}{Laser printers use toner. Ink ribbons are common in impact printers, liquid ink is used in inkjets, and thermal paper is used by thermal printers.}
\JKSSBItem{[Original, MSQ drill]}{\textbf{(MSQ: select ALL correct options.)} Which statements are correct?}
  {\item Dot-matrix printing is impact printing.
   \item Laser printing is non-impact printing.
   \item Inkjet printers place ink droplets on media.
   \item A monitor is a hard-copy device.}
  {A, B, C}{The first three describe printer technologies. A monitor produces soft-copy output, not hard copy.}
\fi
\ifnum\JKSSBLessonID=16
\JKSSBItem{[Original]}{Which connector is commonly used to carry digital video and audio from a computer to a display?}
  {\item HDMI
   \item RJ-11
   \item PS/2
   \item VGA only}
  {A}{HDMI commonly carries digital audio and video. RJ-11 is associated with telephone wiring, PS/2 with legacy keyboard/mouse connections, and VGA carries analog video.}
\JKSSBItem{[Original]}{Which interface is commonly used for a wired Ethernet connection?}
  {\item RJ-45
   \item HDMI
   \item 3.5 mm audio jack
   \item VGA}
  {A}{RJ-45 is commonly used for Ethernet. HDMI and VGA serve displays; a 3.5 mm jack serves analog audio.}
\JKSSBItem{[Original]}{Which statement about USB Type-C is correct?}
  {\item Its plug is reversible.
   \item It is an analog-only video connector.
   \item It is used only for telephone lines.
   \item It is a logical network port number.}
  {A}{USB Type-C is a reversible physical connector. It can carry different signals depending on implementation; the other options confuse it with unrelated interfaces.}
\JKSSBItem{[Original]}{Which port number is a logical transport-layer endpoint rather than a physical socket?}
  {\item TCP port 80
   \item HDMI socket
   \item USB-A receptacle
   \item VGA connector}
  {A}{TCP port 80 is a logical endpoint number. The others are physical connectors.}
\JKSSBItem{[Original, MSQ drill]}{\textbf{(MSQ: select ALL correct options.)} Match the interface to a common use.}
  {\item HDMI --- digital audio/video
   \item RJ-45 --- Ethernet
   \item USB --- peripheral/data connection
   \item VGA --- magnetic-ink character reading}
  {A, B, C}{The first three pairings are standard uses. MICR, not VGA, reads magnetic-ink characters.}
\fi
\ifnum\JKSSBLessonID=17
\JKSSBItem{[Original]}{Which term describes the physical components of a computer?}
  {\item Hardware
   \item Software
   \item Firmware
   \item Data}
  {A}{Hardware is the physical equipment. Software is instructions, firmware is low-level software stored in non-volatile memory, and data is processed content.}
\JKSSBItem{[PYQ-adapted: JKSSB FAA 2024, QCRIAS item 668; SRC-16/SRC-19]}{Firmware is typically stored in which kind of memory?}
  {\item Volatile RAM only
   \item Non-volatile memory such as ROM/flash
   \item CPU cache only
   \item A printer buffer only}
  {B}{Firmware must be available at startup and is stored in non-volatile memory, commonly flash/EEPROM. The other listed memories do not serve as its normal persistent store.}
\JKSSBItem{[Original]}{A device driver is best classified as:}
  {\item Software that helps the operating system communicate with hardware
   \item A physical input device
   \item A storage medium
   \item A type of CPU register}
  {A}{A driver is software that controls or communicates with a device. It is not hardware, storage, or a register.}
\JKSSBItem{[Original]}{Which is an example of application software?}
  {\item Word processor
   \item Device driver
   \item Boot firmware
   \item Memory controller hardware}
  {A}{A word processor helps users perform tasks and is application software. A driver and firmware are system-level software; a memory controller is hardware.}
\JKSSBItem{[Original, MSQ drill]}{\textbf{(MSQ: select ALL correct options.)} Which statements correctly classify components?}
  {\item A keyboard is hardware.
   \item A device driver is software.
   \item Firmware is low-level software associated with hardware control.
   \item A spreadsheet file is a CPU register.}
  {A, B, C}{Keyboard, driver, and firmware classifications are correct. A spreadsheet file is data/document content, not a register.}
\fi
\ifnum\JKSSBLessonID=18
\JKSSBItem{[PYQ-adapted: JKSSB Sub-Inspector 2022, QCRIAS item 329; SRC-17/SRC-22]}{Which software is designed to control computer hardware and provide services for applications?}
  {\item System software
   \item Application software
   \item User data
   \item Presentation theme}
  {A}{System software manages hardware and provides services. Application software performs user tasks; data and themes are not software categories for hardware control.}
\JKSSBItem{[Original]}{Which is a utility program?}
  {\item Disk cleanup tool
   \item Word processor
   \item Presentation file
   \item Device cable}
  {A}{Disk cleanup is a system utility. A word processor is application software; a presentation file is data, and a cable is hardware.}
\JKSSBItem{[Original]}{Freeware is primarily classified by its:}
  {\item Price to the user
   \item Source-code license alone
   \item Processor architecture
   \item Storage medium}
  {A}{Freeware is distributed at no monetary cost, but this alone does not establish open-source freedoms. Processor architecture and storage medium are unrelated.}
\JKSSBItem{[Original]}{Open-source software is distinguished primarily by access to its:}
  {\item Source code under a license that permits specified use/modification
   \item Hardware schematics in every case
   \item Free support hotline
   \item User's private data}
  {A}{Open-source licenses grant rights concerning source code. They do not guarantee free support, hardware schematics, or access to user data.}
\JKSSBItem{[Original, MSQ drill]}{\textbf{(MSQ: select ALL correct options.)} Which are software categories or licensing/distribution labels?}
  {\item System software
   \item Application software
   \item Utility software
   \item Open source}
  {A, B, C, D}{System/application/utility classify roles; open source describes source availability and license terms.}
\fi
\ifnum\JKSSBLessonID=19
\JKSSBItem{[PYQ-adapted: JKSSB Website Operator 2026 Q61, provisional key; SRC-04/SRC-07]}{Which software layer provides common services and an interface between users/applications and computer hardware?}
  {\item Operating system
   \item Spreadsheet
   \item Web page
   \item Device cable}
  {A}{The operating system manages resources and provides services to applications/users. The other options are an application, content, or hardware. The cited key is provisional.}
\JKSSBItem{[PYQ-adapted: JKSSB Junior Assistant 2026 Q70, final key; SRC-03/SRC-06]}{Which action is a recognized way an operating system can recover from deadlock?}
  {\item Resource preemption
   \item Increasing screen brightness
   \item Formatting every disk automatically
   \item Changing a document font}
  {A}{Resource preemption can reclaim resources to break a deadlock. The other actions do not resolve process-resource deadlock.}
\JKSSBItem{[Original]}{Which OS function organizes files into named folders and manages their storage?}
  {\item File-system management
   \item Slide animation
   \item OCR
   \item Mail merge}
  {A}{File-system management organizes stored files/directories. The other options belong to presentation, text recognition, and document automation.}
\JKSSBItem{[Original]}{In a multitasking operating system, the CPU can:}
  {\item Switch among several tasks so they make progress
   \item Execute every program on a separate physical CPU
   \item Store all programs permanently in registers
   \item Eliminate the need for memory}
  {A}{Multitasking schedules CPU time among tasks. It does not require one physical CPU per program or eliminate storage/memory.}
\JKSSBItem{[Original, MSQ drill]}{\textbf{(MSQ: select ALL correct options.)} Which are common operating-system responsibilities?}
  {\item Process scheduling
   \item Memory management
   \item File management
   \item Editing a user's video as its only purpose}
  {A, B, C}{Scheduling processes, managing memory, and managing files are OS responsibilities. Video editing is an application task.}
\fi
\ifnum\JKSSBLessonID=20
\JKSSBItem{[Original]}{Which path format identifies a file by its drive and folders in Windows?}
  {\item C:\textbackslash Users\textbackslash Asha\textbackslash report.docx
   \item https://example.com
   \item user@example.com
   \item 192.0.2.10}
  {A}{A drive-letter path with folders identifies a Windows file location. The other options are a URL, email address, and IP address.}
\JKSSBItem{[Original]}{A file extension usually appears:}
  {\item After the final period in a filename
   \item Before the drive letter
   \item In place of the file contents
   \item Only inside the Recycle Bin}
  {A}{An extension is the filename suffix after the final period, such as `.docx`. It does not replace contents or require the Recycle Bin.}
\JKSSBItem{[Original]}{Which Windows action normally moves a deleted local file to the Recycle Bin rather than immediately erasing it?}
  {\item Delete
   \item Rename
   \item Refresh
   \item Search}
  {A}{Delete commonly moves a file to the Recycle Bin. Rename changes its name, Refresh updates a view, and Search locates items.}
\JKSSBItem{[Original]}{Which operation creates another copy while leaving the original file in place?}
  {\item Copy
   \item Move
   \item Cut
   \item Delete}
  {A}{Copy creates a duplicate. Move relocates the original; Cut prepares it for moving; Delete removes it from its current location.}
\JKSSBItem{[Original, MSQ drill]}{\textbf{(MSQ: select ALL correct options.)} Which statements about files/folders are correct?}
  {\item A folder can contain files and other folders.
   \item A path identifies a location in a file system.
   \item Renaming normally changes a file's name, not its stored contents.
   \item Copying a file necessarily deletes the original.}
  {A, B, C}{Folders can contain items, paths locate them, and renaming changes the name. Copying leaves the original in place.}
\fi
\ifnum\JKSSBLessonID=21
\JKSSBItem{[Original]}{Which shortcut commonly saves the active document in Windows applications?}
  {\item Ctrl+S
   \item Ctrl+P
   \item Ctrl+Z
   \item Ctrl+F}
  {A}{Ctrl+S saves. Ctrl+P prints, Ctrl+Z undoes, and Ctrl+F searches.}
\JKSSBItem{[Original]}{Which shortcut commonly undoes the most recent action?}
  {\item Ctrl+Z
   \item Ctrl+Y
   \item Ctrl+N
   \item Ctrl+O}
  {A}{Ctrl+Z is Undo. Ctrl+Y commonly redoes, Ctrl+N creates a new item, and Ctrl+O opens a file.}
\JKSSBItem{[Original]}{In many Windows applications, Alt+Tab is used to:}
  {\item Switch between open applications
   \item Save the current document
   \item Insert a hyperlink
   \item Close the computer}
  {A}{Alt+Tab switches among open windows. The other actions use different commands or controls.}
\JKSSBItem{[Original]}{Which shortcut usually selects all text or items in the active document/list?}
  {\item Ctrl+A
   \item Ctrl+X
   \item Ctrl+V
   \item Ctrl+K}
  {A}{Ctrl+A selects all. Ctrl+X cuts, Ctrl+V pastes, and Ctrl+K commonly inserts a hyperlink in Office applications.}
\JKSSBItem{[Original, MSQ drill]}{\textbf{(MSQ: select ALL correct options.)} Match the common shortcut to its action.}
  {\item Ctrl+C --- copy
   \item Ctrl+V --- paste
   \item Ctrl+X --- cut
   \item Ctrl+P --- save}
  {A, B, C}{Ctrl+C copies, Ctrl+V pastes, and Ctrl+X cuts. Ctrl+P opens Print; Ctrl+S saves.}
\fi
\ifnum\JKSSBLessonID=22
\JKSSBItem{[PYQ-adapted: SSC CGL 2025 Tier-II Paper-I, 19-Jan-2026 Q14; SRC-18]}
  {A report starts a new section on the next page. Which setting can be changed
  for that section without changing the previous section?}
  {\item AutoCorrect replacement list
   \item Page orientation and margins
   \item Normal style's font face
   \item Document Properties metadata}
  {B}{Page setup belongs to each section, so orientation and margins can differ.
  AutoCorrect is application-wide, the Normal style applies through the document,
  and Properties describe the file rather than one section.}
\JKSSBItem{[PYQ-adapted: JKSSB Junior Assistant 2026 Q73, key:final; SRC-03/SRC-06]}
  {Which break is needed when a later part of a Word document must use a different
  header or footer?}
  {\item Section Break
   \item Page Break
   \item Line Break
   \item Column Break}
  {A}{A Section Break creates a new section with independent header/footer
  settings. A Page Break only starts a page; line and column breaks change text
  flow without creating a section.}
\JKSSBItem{[PYQ-adapted: JKSSB Website Operator 2026 Q77, key:provisional; SRC-04/SRC-07]}
  {Which Word feature combines a main document with a recipient list to produce
  personalized letters?}
  {\item Track Changes
   \item Mail Merge
   \item SmartArt
   \item Word Count}
  {B}{Mail Merge combines a main document with a recipient data source. Track
  Changes records edits, SmartArt creates diagrams, and Word Count reports text
  statistics. The cited key is provisional.}
\JKSSBItem{[Original]}
  {Which shortcut opens the Find and Replace dialog in Word?}
  {\item Ctrl+H
   \item Ctrl+P
   \item Ctrl+B
   \item Ctrl+K}
  {A}{Ctrl+H opens Find and Replace. Ctrl+P opens Print, Ctrl+B toggles bold, and
  Ctrl+K inserts a hyperlink.}
\JKSSBItem{[Original, MSQ drill]}
  {\textbf{(MSQ: select ALL correct options.)} Which settings can be specific to
  a Word section?}
  {\item Page orientation
   \item Page margins
   \item Headers and footers
   \item The font definition of the Normal style}
  {A, B, C}{Orientation, margins, and headers/footers can vary by section. The
  Normal style is a document-level style, not a per-section setting.}
\fi
\ifnum\JKSSBLessonID=23
\JKSSBItem{[PYQ-adapted: SSC CGL 2025 Tier-II Paper-I, 19-Jan-2026 Q16; SRC-18]}{A user changes one heading's appearance but wants other headings using the same style to remain unchanged. What is the safest approach?}
  {\item Modify only that instance's direct formatting
   \item Redefine the shared style
   \item Change the document template for every file
   \item Use AutoCorrect}
  {A}{Direct formatting changes only the selected instance. Redefining a shared style changes all text using it; templates and AutoCorrect do not isolate one heading.}
\JKSSBItem{[Original]}{Which shortcut commonly opens the Find and Replace dialog in Word?}
  {\item Ctrl+H
   \item Ctrl+P
   \item Ctrl+S
   \item Ctrl+L}
  {A}{Ctrl+H opens Find and Replace. Ctrl+P prints, Ctrl+S saves, and Ctrl+L commonly aligns text left.}
\JKSSBItem{[Original]}{Which key commonly starts spelling and grammar checking in Word?}
  {\item F1
   \item F5
   \item F7
   \item F12}
  {C}{F7 starts spelling and grammar review in common Word versions. F1 opens help, F5 navigates/goes to, and F12 opens Save As.}
\JKSSBItem{[Original]}{Which feature records insertions and deletions so a reviewer can accept or reject them?}
  {\item Track Changes
   \item Word Count
   \item Mail Merge
   \item Page Color}
  {A}{Track Changes records edits for review. The other features count words, create personalized documents, or change page appearance.}
\JKSSBItem{[Original, MSQ drill]}{\textbf{(MSQ: select ALL correct options.)} Which features support document review/collaboration?}
  {\item Comments
   \item Track Changes
   \item Accept/Reject Changes
   \item Mail Merge}
  {A, B, C}{Comments and tracked changes support review; Accept/Reject resolves tracked edits. Mail Merge generates personalized documents.}
\fi
\ifnum\JKSSBLessonID=24
\JKSSBItem{[Original]}{The intersection of a worksheet row and column is a:}
  {\item Cell
   \item Workbook
   \item Formula bar
   \item Ribbon}
  {A}{A cell is the intersection of a row and column. A workbook contains worksheets; the formula bar and Ribbon are interface elements.}
\JKSSBItem{[Original]}{A collection of worksheets saved in one Excel file is a:}
  {\item Workbook
   \item Cell range
   \item Chart area
   \item Field}
  {A}{A workbook is the Excel file that can contain worksheets. A range is a set of cells; chart area and field have different meanings.}
\JKSSBItem{[Original]}{Which box displays the address of the active cell and can be used to jump to a cell reference?}
  {\item Name Box
   \item Status bar
   \item Formula bar
   \item Title bar}
  {A}{The Name Box displays the active cell reference and accepts a destination. The Formula bar displays cell content; the status and title bars report other information.}
\JKSSBItem{[Original]}{In the reference B7, which part identifies the column?}
  {\item B
   \item 7
   \item B7
   \item Neither part}
  {A}{The letter identifies the column; the number identifies the row.}
\JKSSBItem{[Original, MSQ drill]}{\textbf{(MSQ: select ALL correct options.)} Which Excel terms are matched correctly?}
  {\item Workbook --- file containing worksheets
   \item Worksheet --- grid of rows and columns
   \item Cell --- intersection of a row and column
   \item Formula bar --- identifies only the active worksheet tab}
  {A, B, C}{The first three definitions are correct. The Formula bar displays or edits the active cell's contents; worksheet tabs identify sheets.}
\fi
\ifnum\JKSSBLessonID=25
\JKSSBItem{[Original]}{Which symbol makes both the row and column absolute in the reference \$A\$1?}
  {\item Dollar sign (\$)
   \item Percent sign (\%)
   \item Hash sign (\#)
   \item Ampersand (\&)}
  {A}{The dollar signs lock the column and row when a formula is copied. The other symbols do not create absolute references.}
\JKSSBItem{[Original]}{When the formula =A1 is copied one column to the right, a relative reference normally becomes:}
  {\item =B1
   \item =\$A\$1
   \item =A\$1
   \item =A2}
  {A}{A relative reference shifts with the copy location, so the column changes from A to B. The other options lock a dimension or change the row.}
\JKSSBItem{[Original]}{Which reference locks the row but allows the column to change?}
  {\item A\$1
   \item \$A1
   \item \$A\$1
   \item A1}
  {A}{A\$1 fixes row 1 but leaves the column relative. \$A1 locks the column; \$A\$1 locks both; A1 locks neither.}
\JKSSBItem{[Original]}{Which operator combines text values in many Excel formulas?}
  {\item \&
   \item \^{}
   \item /
   \item >}
  {A}{The ampersand concatenates text. The other symbols represent exponentiation, division, and comparison.}
\JKSSBItem{[Original, MSQ drill]}{\textbf{(MSQ: select ALL correct options.)} Which are valid Excel cell references?}
  {\item A1
   \item \$B\$7
   \item C\$4
   \item \$D8}
  {A, B, C, D}{All four are valid A1-style references; dollar signs indicate absolute row/column components.}
\fi
\ifnum\JKSSBLessonID=26
\JKSSBItem{[PYQ-adapted: SSC CGL 2025 Tier-II Paper-I, 19-Jan-2026 Q4; SRC-18]}
  {Which Excel function counts cells that satisfy one specified condition?}
  {\item COUNT
   \item COUNTBLANK
   \item COUNTIF
   \item SUMIF}
  {C}{COUNTIF counts cells meeting a criterion. COUNT counts numeric cells
  without a criterion, COUNTBLANK counts empty cells, and SUMIF adds matching
  values rather than counting them.}
\JKSSBItem{[PYQ-adapted: SSC CGL 2025 Tier-II Paper-I, 19-Jan-2026 Q9; SRC-18]}
  {A teacher needs to show only rows where marks exceed 75 while keeping the
  remaining records in the worksheet. Which tool should be used?}
  {\item Sort
   \item Filter
   \item Remove Duplicates
   \item Subtotal}
  {B}{Filter temporarily displays rows that meet criteria. Sort changes order,
  Remove Duplicates deletes repeated records, and Subtotal inserts grouped
  summaries.}
\JKSSBItem{[Original]}
  {Which function counts non-empty cells, including cells containing text?}
  {\item COUNTA
   \item COUNT
   \item COUNTBLANK
   \item SUM}
  {A}{COUNTA counts non-empty cells. COUNT counts numeric cells, COUNTBLANK
  counts empty cells, and SUM adds numeric values.}
\JKSSBItem{[Original]}
  {In a formula, which function tests a condition and returns one value when it
  is true and another when it is false?}
  {\item MAX
   \item IF
   \item AVERAGE
   \item ROUND}
  {B}{IF evaluates a logical test and selects between two results. MAX returns
  the largest value, AVERAGE calculates a mean, and ROUND changes numeric
  precision.}
\JKSSBItem{[Original, MSQ drill]}
  {\textbf{(MSQ: select ALL correct options.)} Which statements about Excel
  counting functions are correct?}
  {\item COUNT counts cells containing numbers.
   \item COUNTA counts non-empty cells.
   \item COUNTBLANK counts empty cells.
   \item COUNTIF sums the values that meet a criterion.}
  {A, B, C}{COUNT, COUNTA, and COUNTBLANK perform the stated counts. COUNTIF
  counts matching cells; SUMIF, not COUNTIF, sums their values.}
\fi
\ifnum\JKSSBLessonID=27
\JKSSBItem{[PYQ-adapted: SSC CGL 2025 Tier-II Paper-I, 19-Jan-2026 Q9; SRC-18]}{A teacher needs to show only rows where marks exceed 75 while keeping other records. Which tool should be used?}
  {\item Filter
   \item Remove Duplicates
   \item Delete Sheet
   \item Sort A to Z}
  {A}{Filter displays rows meeting criteria without deleting the rest. Remove Duplicates removes records; deleting a sheet removes content; sorting changes order.}
\JKSSBItem{[Original]}{Which Excel feature summarizes and reorganizes a large data table by categories?}
  {\item PivotTable
   \item AutoCorrect
   \item Mail Merge
   \item Slide Master}
  {A}{A PivotTable summarizes and rearranges data by fields. The other features belong to text correction, Word, or PowerPoint.}
\JKSSBItem{[Original]}{Which chart is generally suitable for showing a trend over time?}
  {\item Line chart
   \item Pie chart
   \item Doughnut chart
   \item Radar chart}
  {A}{A line chart makes change over an ordered time axis easy to see. Pie/doughnut charts emphasize parts of a whole; radar charts compare dimensions.}
\JKSSBItem{[Original]}{Which tool adds drop-down controls to choose which records are displayed?}
  {\item AutoFilter
   \item Goal Seek
   \item Spell Check
   \item Page Break Preview}
  {A}{AutoFilter provides criteria drop-downs. Goal Seek solves for an input; the other tools serve unrelated functions.}
\JKSSBItem{[Original, MSQ drill]}{\textbf{(MSQ: select ALL correct options.)} Which statements about Excel data tools are correct?}
  {\item Sorting changes record order.
   \item Filtering hides non-matching rows from the current view.
   \item A PivotTable can aggregate data by fields.
   \item Removing duplicates only changes font formatting.}
  {A, B, C}{Sorting changes order, filtering changes visible rows, and PivotTables aggregate. Remove Duplicates changes records, not just formatting.}
\fi
\ifnum\JKSSBLessonID=28
\JKSSBItem{[Original]}{In a relational database table, a row is commonly called a:}
  {\item Record
   \item Field
   \item Query
   \item Form}
  {A}{A row is a record; a column is a field. Queries retrieve/process data and forms provide an interface.}
\JKSSBItem{[Original]}{In MS Access, a field corresponds most closely to a table:}
  {\item Column
   \item Row
   \item Report
   \item Database file}
  {A}{A field is a column storing one attribute. A row is a record; reports and database files are higher-level objects.}
\JKSSBItem{[Original]}{Which data type is most appropriate for a field storing a calendar date?}
  {\item Date/Time
   \item Currency
   \item Attachment
   \item Yes/No}
  {A}{Date/Time stores date and time values. Currency stores money; Attachment stores files; Yes/No stores Boolean values.}
\JKSSBItem{[Original]}{A DBMS primarily helps users:}
  {\item Define, store, retrieve, and manage data
   \item Design computer circuits
   \item Print only slide presentations
   \item Replace all operating systems}
  {A}{A DBMS manages databases and access to their data. The other options describe unrelated tasks.}
\JKSSBItem{[Original, MSQ drill]}{\textbf{(MSQ: select ALL correct options.)} Which are common database objects in Access?}
  {\item Tables
   \item Queries
   \item Forms
   \item Reports}
  {A, B, C, D}{Tables store data, queries retrieve/process it, forms support entry/viewing, and reports format output.}
\fi
\ifnum\JKSSBLessonID=29
\JKSSBItem{[Original]}{A primary key should identify each record in a table:}
  {\item Uniquely
   \item With a duplicate value in every row
   \item Only by display color
   \item Using a blank value}
  {A}{A primary key uniquely identifies records and cannot be null. The other choices defeat reliable identification.}
\JKSSBItem{[Original]}{A foreign key is used primarily to:}
  {\item Reference a related record in another table
   \item Encrypt the database file
   \item Sort a report alphabetically
   \item Store an image only}
  {A}{A foreign key links a record to a key in another table and supports referential integrity. The alternatives are unrelated.}
\JKSSBItem{[Original]}{A key made from two or more fields is called a:}
  {\item Composite key
   \item Foreign form
   \item Query key
   \item Null key}
  {A}{A composite key uses multiple fields together to identify a record. The other terms are not standard key categories.}
\JKSSBItem{[Original]}{Referential integrity helps ensure that:}
  {\item Relationships between tables remain valid
   \item Every field is printed in bold
   \item Passwords are stored as plain text
   \item A query has no criteria}
  {A}{Referential integrity constrains relationships, such as preventing unmatched foreign-key references. The other options are unrelated.}
\JKSSBItem{[Original, MSQ drill]}{\textbf{(MSQ: select ALL correct options.)} Which statements about keys are correct?}
  {\item A primary key uniquely identifies a row.
   \item A foreign key can reference a primary key in another table.
   \item A composite key contains multiple fields.
   \item A primary key should contain duplicate values.}
  {A, B, C}{The first three statements are correct. Duplicate primary-key values make records ambiguous.}
\fi
\ifnum\JKSSBLessonID=30
\JKSSBItem{[Original]}{Which Access object retrieves selected records based on criteria?}
  {\item Query
   \item Form
   \item Report
   \item Macro}
  {A}{A query selects or manipulates data based on criteria. Forms support interaction, reports format output, and macros automate actions.}
\JKSSBItem{[Original]}{Which object is primarily designed for a formatted, printable summary of data?}
  {\item Report
   \item Table
   \item Relationship
   \item Module}
  {A}{A report formats data for viewing or printing. Tables store records; relationships connect tables; modules contain code.}
\JKSSBItem{[Original]}{Which Access object is commonly used as a user-friendly screen for entering records?}
  {\item Form
   \item Report
   \item Index
   \item Relationship}
  {A}{Forms provide an interface to enter or view data. Reports present formatted output; indexes and relationships support data organization.}
\JKSSBItem{[Original]}{A SELECT query is generally used to:}
  {\item Retrieve records meeting conditions
   \item Delete the database application
   \item Change the computer's BIOS
   \item Print a PowerPoint slide}
  {A}{A SELECT query returns records from one or more tables based on criteria. The other choices are unrelated.}
\JKSSBItem{[Original, MSQ drill]}{\textbf{(MSQ: select ALL correct options.)} Match each Access object to its typical role.}
  {\item Table --- store records
   \item Query --- retrieve selected data
   \item Form --- data entry/view interface
   \item Report --- formatted presentation of data}
  {A, B, C, D}{All four are standard roles of Access objects.}
\fi
\ifnum\JKSSBLessonID=31
\JKSSBItem{[Original]}{A PowerPoint slide layout controls the:}
  {\item Arrangement of placeholders on a slide
   \item Order of files in a folder
   \item Worksheet cell address
   \item Email server protocol}
  {A}{A layout arranges slide placeholders. The other options concern file management, Excel, or email.}
\JKSSBItem{[Original]}{Which PowerPoint feature stores common formatting and repeated slide elements?}
  {\item Slide Master
   \item Formula bar
   \item Navigation Pane
   \item Mail Merge}
  {A}{Slide Master controls shared design and layout elements. The alternatives are Excel/Word features.}
\JKSSBItem{[Original]}{A placeholder on a slide is intended to:}
  {\item Hold content such as text, pictures, or charts
   \item Store an operating system
   \item Save an email password
   \item Execute a database query}
  {A}{Placeholders reserve positions for slide content. They do not store an OS, credentials, or run database queries.}
\JKSSBItem{[Original]}{Which view is designed to rearrange slide thumbnails quickly?}
  {\item Slide Sorter
   \item Notes Page
   \item Reading View
   \item Outline view}
  {A}{Slide Sorter presents thumbnails for rearranging slides. Notes Page emphasizes speaker notes; Reading View presents the show; Outline view focuses on text.}
\JKSSBItem{[Original, MSQ drill]}{\textbf{(MSQ: select ALL correct options.)} Which are common PowerPoint views?}
  {\item Normal view
   \item Slide Sorter view
   \item Notes Page view
   \item Slide Master view}
  {A, B, C, D}{All four are PowerPoint views used for editing, arranging, notes, or shared design.}
\fi
\ifnum\JKSSBLessonID=32
\JKSSBItem{[Original]}{A transition effect occurs when:}
  {\item Moving from one slide to another
   \item A word is selected in a text box
   \item A spreadsheet is recalculated
   \item A file is renamed}
  {A}{Transitions apply between slides. Animations apply to objects within a slide; the other actions are unrelated.}
\JKSSBItem{[PYQ-adapted: JKSSB Website Operator 2026 Q80, provisional key; SRC-04/SRC-07]}{Which key starts a PowerPoint slide show from the beginning?}
  {\item F5
   \item Shift+F5
   \item F7
   \item Ctrl+S}
  {A}{F5 starts from the beginning. Shift+F5 starts from the current slide, F7 commonly checks spelling, and Ctrl+S saves. The cited key is provisional.}
\JKSSBItem{[Original]}{Which shortcut commonly starts a slide show from the current slide?}
  {\item Shift+F5
   \item F5
   \item Ctrl+P
   \item Alt+F4}
  {A}{Shift+F5 starts at the current slide. F5 starts at the beginning; Ctrl+P prints; Alt+F4 closes the active window.}
\JKSSBItem{[Original]}{An animation effect is applied to:}
  {\item An object or content within a slide
   \item The change between two files in a folder
   \item A worksheet's row height only
   \item The computer boot process}
  {A}{Animations affect objects on a slide; transitions govern movement between slides. The other choices are unrelated.}
\JKSSBItem{[Original, MSQ drill]}{\textbf{(MSQ: select ALL correct options.)} Which statements are correct?}
  {\item F5 starts the slide show from the beginning.
   \item Shift+F5 starts from the current slide.
   \item A transition is between slides.
   \item An animation is always the same as a transition.}
  {A, B, C}{The first three describe standard functions. Animation and transition refer to different scopes.}
\fi
\ifnum\JKSSBLessonID=33
\JKSSBItem{[Original]}{Which file extension is commonly associated with a modern Word document?}
  {\item .docx
   \item .xlsx
   \item .pptx
   \item .accdb}
  {A}{.docx is Word, .xlsx is Excel, .pptx is PowerPoint, and .accdb is Access.}
\JKSSBItem{[Original]}{A PDF is primarily designed to preserve:}
  {\item Document layout across devices
   \item Spreadsheet formulas as the only content
   \item Executable program instructions
   \item Audio channels only}
  {A}{PDF is a page-description/document format intended to preserve layout. It is not a spreadsheet-only or executable format.}
\JKSSBItem{[Original]}{Which file format is commonly used for plain comma-separated tabular data?}
  {\item CSV
   \item PPTX
   \item PNG
   \item EXE}
  {A}{CSV stores delimited text rows and columns. PPTX is a presentation package, PNG is an image format, and EXE is an executable extension.}
\JKSSBItem{[Original]}{Which format is an image format?}
  {\item PNG
   \item DOCX
   \item XLSX
   \item PPTX}
  {A}{PNG is a raster image format. DOCX, XLSX, and PPTX are Office document formats.}
\JKSSBItem{[Original, MSQ drill]}{\textbf{(MSQ: select ALL correct options.)} Which extension-to-file-type pairings are correct?}
  {\item .docx --- word-processing document
   \item .xlsx --- spreadsheet workbook
   \item .pptx --- presentation
   \item .png --- database file}
  {A, B, C}{The first three pairings are correct. PNG is an image format, not a database file.}
\fi
\ifnum\JKSSBLessonID=34
\JKSSBItem{[Original]}{The World Wide Web is best described as:}
  {\item A system of interlinked resources accessed over the Internet
   \item The physical cables of the Internet
   \item A single search engine
   \item A computer's operating system}
  {A}{The Web is an information system of linked resources that uses the Internet. It is not the network infrastructure, one search engine, or an OS.}
\JKSSBItem{[Original]}{Which application retrieves and displays web pages?}
  {\item Web browser
   \item Search engine
   \item Spreadsheet
   \item Mail server}
  {A}{A browser requests and displays web resources. A search engine helps locate resources; the other choices serve different roles.}
\JKSSBItem{[Original]}{A search engine primarily helps users:}
  {\item Find indexed web content
   \item Route packets at the network layer
   \item Edit local spreadsheets
   \item Store BIOS firmware}
  {A}{Search engines index and retrieve web information. Routing is a network function; spreadsheets and firmware storage are unrelated.}
\JKSSBItem{[Original]}{A single document available on the Web is a:}
  {\item Web page
   \item Website
   \item Browser
   \item Domain registrar}
  {A}{A web page is an individual resource/document. A website is a collection of related pages; the other options are software/services.}
\JKSSBItem{[Original, MSQ drill]}{\textbf{(MSQ: select ALL correct options.)} Which statements are correct?}
  {\item The Internet is a global network of networks.
   \item The Web uses Internet infrastructure.
   \item A browser and a search engine are different tools.
   \item Every Internet service is a web page.}
  {A, B, C}{The Internet supports many services; the Web is one of them. Browsers and search engines also have distinct roles.}
\fi
\ifnum\JKSSBLessonID=35
\JKSSBItem{[Original]}{In a URL, the scheme identifies the:}
  {\item Access method or protocol
   \item User's computer model
   \item Page's font size
   \item Local file extension only}
  {A}{The scheme, such as https, indicates the access protocol. The other details are not what a URL scheme encodes.}
\JKSSBItem{[Original]}{Which service translates a domain name into an IP address?}
  {\item DNS
   \item DHCP
   \item SMTP
   \item FTP}
  {A}{DNS resolves names to records such as IP addresses. DHCP assigns network configuration; SMTP sends mail; FTP transfers files.}
\JKSSBItem{[Original]}{HTTPS adds which security protocol to web communication?}
  {\item TLS
   \item POP3
   \item OMR
   \item RAID}
  {A}{HTTPS uses HTTP protected by TLS. POP3 receives email, OMR reads marks, and RAID organizes disks.}
\JKSSBItem{[Original]}{Which component of the URL `example.org` is the top-level domain?}
  {\item .org
   \item example
   \item https
   \item /}
  {A}{`.org` is the top-level domain. `example` is the domain label; https is a scheme, and slash separates path components.}
\JKSSBItem{[Original, MSQ drill]}{\textbf{(MSQ: select ALL correct options.)} Which statements about web addressing/security are correct?}
  {\item DNS can resolve domain names to IP addresses.
   \item HTTPS uses TLS to protect a connection.
   \item A URL can include a scheme and host.
   \item DNS assigns an email password to a user.}
  {A, B, C}{DNS resolves names, HTTPS uses TLS, and URLs include components such as scheme and host. DNS does not manage passwords.}
\fi
\ifnum\JKSSBLessonID=36
\JKSSBItem{[Original]}{Downloading a file means transferring it:}
  {\item From a remote system to the local device
   \item From the local device to a remote system
   \item From RAM to the CPU register file only
   \item From a monitor to a printer}
  {A}{Download means receiving data from a remote service. Sending data from the local device is upload.}
\JKSSBItem{[Original]}{A website crawler primarily:}
  {\item Discovers and visits web pages for indexing
   \item Encrypts every email attachment
   \item Replaces the browser's operating system
   \item Assigns IP addresses on a LAN}
  {A}{A crawler discovers pages that may be indexed by a search engine. The other options describe unrelated security, OS, or network services.}
\JKSSBItem{[Original]}{Browser cache is used mainly to:}
  {\item Temporarily store resources for faster reuse
   \item Permanently back up a user's documents
   \item Replace DNS
   \item Prevent all malware infections}
  {A}{Cache can speed repeat access by storing copies temporarily. It is not a permanent backup, DNS replacement, or complete security tool.}
\JKSSBItem{[Original]}{Which action sends a local file to a website or remote service?}
  {\item Upload
   \item Download
   \item Refresh
   \item Bookmark}
  {A}{Upload sends data from local to remote. Download receives remote data; refresh reloads a page; bookmark saves a link.}
\JKSSBItem{[Original, MSQ drill]}{\textbf{(MSQ: select ALL correct options.)} Which statements are correct?}
  {\item Browser history can record visited pages.
   \item Cookies can store site-related data in a browser.
   \item Upload sends data to a remote service.
   \item Browser cache is a guaranteed disaster-recovery backup.}
  {A, B, C}{History records visits, cookies store site data, and upload sends data outward. Cache is temporary and is not a guaranteed backup.}
\fi
\ifnum\JKSSBLessonID=37
\JKSSBItem{[Original]}{Which email field hides listed recipients from other recipients?}
  {\item Bcc
   \item To
   \item Cc
   \item Subject}
  {A}{Bcc hides those recipients from other recipients. To and Cc recipients are visible; Subject is the message topic.}
\JKSSBItem{[Original]}{A file sent along with an email is called an:}
  {\item Attachment
   \item Header
   \item Signature
   \item Folder}
  {A}{An attachment is a file included with a message. Headers carry routing/message metadata; a signature identifies the sender; a folder organizes mail.}
\JKSSBItem{[Original]}{Which email action sends a received message to another recipient?}
  {\item Forward
   \item Reply
   \item Draft
   \item Archive}
  {A}{Forward sends a received message onward. Reply responds to its sender; Draft stores an unfinished message; Archive files a message.}
\JKSSBItem{[Original]}{In an email address, the `@` symbol separates the:}
  {\item User name and domain name
   \item Scheme and file path
   \item Subject and message body
   \item Password and server port}
  {A}{The local part identifies the mailbox/user; the domain identifies the mail domain. The other pairs are unrelated.}
\JKSSBItem{[Original, MSQ drill]}{\textbf{(MSQ: select ALL correct options.)} Which statements about email fields are correct?}
  {\item To lists primary recipients.
   \item Cc lists visible copied recipients.
   \item Bcc conceals listed recipients from other recipients.
   \item Subject contains the attachment bytes by definition.}
  {A, B, C}{To, Cc, and Bcc serve recipient roles. Subject is a text topic field; attachments are separate message parts.}
\fi
\ifnum\JKSSBLessonID=38
\JKSSBItem{[PYQ-adapted: JKSSB Sub-Inspector 2022, QCRIAS item 399; SRC-17/SRC-22]}{Which protocol is commonly used by an email client to retrieve messages from a remote server?}
  {\item SMTP
   \item POP3
   \item SNMP
   \item UDP}
  {B}{POP3 retrieves mail from a server. SMTP sends/transfers email; SNMP manages network devices; UDP is a transport protocol.}
\JKSSBItem{[Original]}{Which protocol is primarily used to send email from a client to a mail server?}
  {\item SMTP
   \item IMAP
   \item POP3
   \item DNS}
  {A}{SMTP handles mail submission/transfer. IMAP and POP3 retrieve/access mail; DNS resolves names.}
\JKSSBItem{[Original]}{Which protocol typically keeps messages synchronized between a client and a mail server?}
  {\item IMAP
   \item POP3 only
   \item FTP
   \item ARP}
  {A}{IMAP is designed for server-based mailbox access and synchronization. POP3 commonly downloads messages; FTP transfers files; ARP resolves local addresses.}
\JKSSBItem{[Original]}{TLS applied to SMTP is intended to protect:}
  {\item The transport connection between mail systems
   \item The physical paper envelope
   \item The computer's RAM contents at shutdown
   \item Printer output resolution}
  {A}{TLS can protect a connection hop. It does not by itself provide end-to-end message encryption across every mail server.}
\JKSSBItem{[Original, MSQ drill]}{\textbf{(MSQ: select ALL correct options.)} Match the email protocol to its common role.}
  {\item SMTP --- sending/transfer
   \item POP3 --- retrieval/download
   \item IMAP --- server mailbox access/synchronization
   \item DNS --- composing email text}
  {A, B, C}{SMTP, POP3, and IMAP have the listed mail roles. DNS resolves domain names, not message composition.}
\fi
\ifnum\JKSSBLessonID=39
\JKSSBItem{[Original]}{An OTP is best described as:}
  {\item A one-time verification code
   \item A permanent account password
   \item A bank account number
   \item A card's magnetic strip}
  {A}{An OTP is a temporary one-time code for authentication. It is not a permanent password or payment identifier.}
\JKSSBItem{[Original]}{Which detail should never be shared with an unsolicited caller asking to authorize a payment?}
  {\item UPI PIN
   \item Public merchant name
   \item Payment date
   \item Transaction category}
  {A}{A UPI PIN authorizes transactions and must remain secret. Public transaction context does not substitute for an authentication secret.}
\JKSSBItem{[Original]}{A collect request in a payment app should be treated as:}
  {\item A request to approve a payment, which must be checked carefully
   \item Proof that money has already arrived
   \item A harmless receipt
   \item A request to disclose an OTP to the sender}
  {A}{Approving a collect request may authorize a debit. Check the payee and purpose; it is not proof of receipt and never requires sharing an OTP.}
\JKSSBItem{[Original]}{Which security practice adds a second verification factor beyond a password?}
  {\item Multi-factor authentication
   \item Reusing the same password
   \item Disabling screen lock
   \item Saving a PIN in a public note}
  {A}{MFA combines distinct factors. The other choices weaken security.}
\JKSSBItem{[Original, MSQ drill]}{\textbf{(MSQ: select ALL correct options.)} Which are safe digital-payment practices?}
  {\item Verify the recipient before confirming.
   \item Never disclose a UPI PIN or OTP to another person.
   \item Reject unexpected collect requests.
   \item Enter a PIN to receive money into an account.}
  {A, B, C}{Verify recipient details and protect secrets. Receiving money does not require disclosing or entering a UPI PIN to a requester.}
\fi
\ifnum\JKSSBLessonID=40
\JKSSBItem{[PYQ-adapted: JKSSB FAA 2022, QCRIAS item 547; SRC-21]}{The Internet is an example of which network type by geographic scope?}
  {\item LAN
   \item MAN
   \item WAN
   \item PAN}
  {C}{The Internet spans large geographic areas and interconnected networks, so it is a WAN. LAN, MAN, and PAN cover smaller scopes.}
\JKSSBItem{[PYQ-adapted: JKSSB Sub-Inspector 2022, QCRIAS item 408; SRC-17/SRC-22]}{Which device directs traffic between networks?}
  {\item Router
   \item Hub
   \item Keyboard
   \item Plotter}
  {A}{A router forwards packets between networks. A hub repeats local signals; keyboard is input and plotter is output.}
\JKSSBItem{[Original]}{Which topology connects each node to a central device?}
  {\item Star
   \item Ring
   \item Bus
   \item Mesh only}
  {A}{In a star topology each node connects to a central hub/switch. Ring and bus use different arrangements; mesh uses multiple inter-node paths.}
\JKSSBItem{[Original]}{A device that connects a computer to a network is commonly a:}
  {\item NIC
   \item Plotter
   \item OCR reader
   \item Projector}
  {A}{A Network Interface Card (NIC) provides network connectivity. The others are output or recognition devices.}
\JKSSBItem{[Original, MSQ drill]}{\textbf{(MSQ: select ALL correct options.)} Which statements are correct?}
  {\item A LAN covers a limited local area.
   \item A router forwards traffic between networks.
   \item A switch commonly connects devices within a LAN.
   \item A hub makes routing decisions using IP addresses.}
  {A, B, C}{LAN, router, and switch roles are correctly described. A hub repeats signals and does not route using IP addresses.}
\fi
\ifnum\JKSSBLessonID=41
\JKSSBItem{[PYQ-adapted: JKSSB JKP Constable 2024, QCRIAS item 747; SRC-02]}{In HTTPS, what does the letter S stand for?}
  {\item Selected
   \item Secure
   \item Software
   \item System}
  {B}{The S in HTTPS means Secure. The other options are not expansions of HTTPS.}
\JKSSBItem{[PYQ-adapted: JKSSB JKP Constable 2024, QCRIAS item 765; SRC-02]}{Which organization develops standards for the World Wide Web?}
  {\item W3C
   \item WHO
   \item OPEC
   \item UNESCO}
  {A}{W3C is the World Wide Web Consortium. The other organizations have different mandates.}
\JKSSBItem{[Original]}{Which protocol provides reliable, ordered delivery between applications?}
  {\item TCP
   \item UDP
   \item ARP
   \item ICMP}
  {A}{TCP provides reliable, ordered transport. UDP is connectionless without that delivery guarantee; ARP and ICMP serve other network functions.}
\JKSSBItem{[Original]}{Which protocol translates a domain name into an IP address?}
  {\item DNS
   \item SMTP
   \item FTP
   \item POP3}
  {A}{DNS resolves names. SMTP sends mail, FTP transfers files, and POP3 retrieves email.}
\JKSSBItem{[Original, MSQ drill]}{\textbf{(MSQ: select ALL correct options.)} Which are application-layer protocols/services?}
  {\item HTTP
   \item SMTP
   \item FTP
   \item Ethernet framing}
  {A, B, C}{HTTP, SMTP, and FTP are application-layer protocols. Ethernet framing belongs to the data-link layer.}
\fi
\ifnum\JKSSBLessonID=42
\JKSSBItem{[PYQ-adapted: JKSSB FAA 2022, QCRIAS item 556; SRC-21]}{Which virus type attaches itself to executable files and activates when an infected file runs?}
  {\item Direct-action file virus
   \item Boot-sector virus
   \item Macro virus
   \item Overwrite virus only}
  {A}{A direct-action file virus infects executable files such as .EXE/.COM files. Boot-sector viruses target boot records; macro viruses use document macros.}
\JKSSBItem{[PYQ-adapted: JKSSB FAA 2024, QCRIAS item 664; SRC-16/SRC-19]}{Which option is an antivirus product?}
  {\item Trojan horse
   \item Worm
   \item Macro virus
   \item Norton}
  {D}{Norton is antivirus software. A Trojan, worm, and macro virus are malware types.}
\JKSSBItem{[Original]}{Malware that encrypts a victim's files and demands payment is commonly called:}
  {\item Ransomware
   \item Adware
   \item Cookie
   \item Device driver}
  {A}{Ransomware encrypts or locks data and demands a ransom. Adware displays advertisements; a cookie is browser data; a driver is software for a device.}
\JKSSBItem{[Original]}{Which malware is commonly disguised as legitimate software and relies on a user running it?}
  {\item Trojan horse
   \item Worm
   \item Firewall
   \item Patch}
  {A}{A Trojan appears benign but carries a hidden malicious purpose. A worm self-propagates; firewalls and patches are defensive controls.}
\JKSSBItem{[Original, MSQ drill]}{\textbf{(MSQ: select ALL correct options.)} Which are malware categories?}
  {\item Virus
   \item Worm
   \item Trojan horse
   \item Firewall}
  {A, B, C}{Virus, worm, and Trojan horse are malware. A firewall is a security control.}
\fi
\ifnum\JKSSBLessonID=43
\JKSSBItem{[Original]}{A fraudulent message that impersonates a trusted organization to steal credentials is an example of:}
  {\item Phishing
   \item Defragmentation
   \item Encryption
   \item Load balancing}
  {A}{Phishing uses deceptive communication to obtain sensitive information. The other terms describe storage, security, or traffic-management processes.}
\JKSSBItem{[Original]}{Forging the apparent sender address of an email is called:}
  {\item Spoofing
   \item Patching
   \item Hashing
   \item Sandboxing}
  {A}{Spoofing falsifies an identity or source. Patching updates software; hashing computes a digest; sandboxing isolates execution.}
\JKSSBItem{[Original]}{An attack that tricks a person into revealing confidential information primarily exploits:}
  {\item Social engineering
   \item Disk formatting
   \item Data compression
   \item Routing}
  {A}{Social engineering manipulates people rather than relying solely on a technical vulnerability. The other options are unrelated operations.}
\JKSSBItem{[Original]}{A software weakness that could be exploited by an attacker is a:}
  {\item Vulnerability
   \item Patch
   \item Backup
   \item Checksum}
  {A}{A vulnerability is a weakness; an exploit takes advantage of it. A patch fixes defects, a backup preserves data, and a checksum detects changes.}
\JKSSBItem{[Original, MSQ drill]}{\textbf{(MSQ: select ALL correct options.)} Which are common signs of phishing?}
  {\item A mismatched sender/domain
   \item An urgent request to disclose a password
   \item A link whose destination differs from its displayed text
   \item A routine operating-system update from the configured updater}
  {A, B, C}{Sender mismatches, credential pressure, and deceptive links are warning signs. A routine configured update is not by itself evidence of phishing.}
\fi
\ifnum\JKSSBLessonID=44
\JKSSBItem{[Original]}{Which security property means that information is accessible only to authorized parties?}
  {\item Confidentiality
   \item Availability
   \item Redundancy
   \item Compression}
  {A}{Confidentiality restricts disclosure to authorized parties. Availability concerns access when needed; redundancy and compression are different concepts.}
\JKSSBItem{[Original]}{A firewall primarily helps to:}
  {\item Control network traffic according to rules
   \item Create a backup copy of every document
   \item Repair broken monitor pixels
   \item Convert scanned pages to editable text}
  {A}{A firewall filters or controls network traffic. Backups, monitor repair, and OCR are separate functions.}
\JKSSBItem{[Original]}{Why should security updates be installed promptly?}
  {\item They can fix known vulnerabilities
   \item They guarantee that no attack can occur
   \item They replace all backups
   \item They make passwords unnecessary}
  {A}{Updates often patch known weaknesses but cannot guarantee complete security or replace backups/authentication.}
\JKSSBItem{[Original]}{Which control is most directly used to restore files after accidental deletion or disk failure?}
  {\item Backup
   \item Firewall
   \item Screen saver
   \item Spam filter}
  {A}{A backup provides a recoverable copy. Firewalls filter traffic, screen savers lock/display, and spam filters classify email.}
\JKSSBItem{[Original, MSQ drill]}{\textbf{(MSQ: select ALL correct options.)} Which practices strengthen account security?}
  {\item Use multi-factor authentication.
   \item Use unique strong passwords.
   \item Install security updates.
   \item Reuse one password on every service.}
  {A, B, C}{MFA, unique passwords, and updates strengthen security. Password reuse increases the impact of a compromised account.}
\fi
\ifnum\JKSSBLessonID=45
\JKSSBItem{[Original]}{Open-source software generally makes which material available under its license?}
  {\item Source code
   \item The user's password database
   \item Private files of all users
   \item A computer's BIOS password}
  {A}{Open-source licenses provide rights to inspect/use/modify source under their terms. They do not expose private user data or credentials.}
\JKSSBItem{[Original]}{Freeware and open-source are:}
  {\item Different concepts; free price does not necessarily grant source-code rights
   \item Always identical terms
   \item Both hardware categories
   \item Names for trial versions only}
  {A}{Freeware concerns price; open source concerns source access and license freedoms. They may overlap but are not synonymous.}
\JKSSBItem{[Original]}{Which license-related statement is most accurate?}
  {\item Users must follow the software's license terms.
   \item Open-source software has no license.
   \item Public-domain software is always proprietary.
   \item Freeware may never be copyrighted.}
  {A}{Open-source and other software are distributed under licenses/terms. The other claims are false generalizations.}
\JKSSBItem{[Original]}{A proprietary application is usually distributed with:}
  {\item Restrictions defined by its owner/license
   \item Automatic public rights to modify its source
   \item No copyright protection
   \item Mandatory source publication}
  {A}{Proprietary licensing typically restricts use/modification according to the owner's terms. Source publication and modification rights are not automatic.}
\JKSSBItem{[Original, MSQ drill]}{\textbf{(MSQ: select ALL correct options.)} Which statements about software distribution are correct?}
  {\item Freeware may be free of charge while remaining closed-source.
   \item Open-source software is governed by a license.
   \item Shareware may be distributed as a trial.
   \item All free-of-charge software is public domain.}
  {A, B, C}{Freeware, open-source, and shareware describe different distribution/licensing properties. No-cost distribution does not automatically place software in the public domain.}
\fi
\ifnum\JKSSBLessonID=46
\JKSSBItem{[PYQ-adapted: JKSSB FAA 2024, QCRIAS item 668; SRC-16/SRC-19]}{Firmware is commonly stored in:}
  {\item Cache
   \item RAM
   \item External printer memory only
   \item ROM or flash memory}
  {D}{Firmware must persist without power and is commonly stored in ROM/flash. Cache and RAM are volatile; printer memory is not the general answer.}
\JKSSBItem{[PYQ-adapted: JKSSB Website Operator 2026 Q68, provisional key; SRC-04/SRC-07]}{Which firmware initializes hardware and begins the boot process on a traditional PC?}
  {\item BIOS
   \item Device driver
   \item Spreadsheet
   \item Web browser}
  {A}{BIOS firmware initializes hardware and starts booting. Drivers are loaded for device operation; spreadsheet/browser are applications. The cited key is provisional.}
\JKSSBItem{[Original]}{UEFI is best described as:}
  {\item A modern firmware interface used during startup
   \item A type of spreadsheet formula
   \item A printer cable
   \item A file-compression format}
  {A}{UEFI is firmware interface technology used to initialize the system and start the boot process. The other choices are unrelated.}
\JKSSBItem{[Original]}{A device driver primarily enables:}
  {\item The operating system to communicate with a hardware device
   \item A monitor to store files permanently
   \item A user to change the CPU clock by typing a document
   \item A network to replace the operating system}
  {A}{Drivers let the OS control/communicate with hardware. The other options confuse storage, applications, and system software.}
\JKSSBItem{[Original, MSQ drill]}{\textbf{(MSQ: select ALL correct options.)} Which statements are correct?}
  {\item Firmware is low-level software associated with a device.
   \item ROM/flash can retain firmware without power.
   \item BIOS/UEFI participates in startup.
   \item A device driver is a physical circuit board.}
  {A, B, C}{Firmware is software stored persistently; BIOS/UEFI participates in boot. A driver is software, not a circuit board.}
\fi
\ifnum\JKSSBLessonID=47
\JKSSBItem{[Original]}{In an information system, the people who use, operate, or manage the system are sometimes collectively called:}
  {\item Humanware
   \item Firmware
   \item Shareware
   \item Middleware}
  {A}{Humanware refers to people involved with computer systems. Firmware, shareware, and middleware are software terms.}
\JKSSBItem{[Original]}{IT governance is primarily concerned with:}
  {\item Ensuring IT supports organizational goals and is responsibly managed
   \item Increasing screen resolution only
   \item Replacing all staff with computers
   \item Removing accountability for systems}
  {A}{IT governance aligns and oversees technology, risk, and accountability in support of organizational objectives. The other choices are too narrow or incorrect.}
\JKSSBItem{[Original]}{Which role is most associated with managing user accounts and access permissions?}
  {\item System administrator
   \item End user only
   \item Printer operator only
   \item Search engine crawler}
  {A}{System administrators commonly manage users, permissions, and systems. The other roles do not generally control organization-wide access.}
\JKSSBItem{[Original]}{An end user is someone who primarily:}
  {\item Uses an application or system to perform a task
   \item Writes every operating-system kernel
   \item Manufactures all processors
   \item Routes every Internet packet}
  {A}{An end user uses the system for work or personal tasks. The other options describe specialized engineering or network functions.}
\JKSSBItem{[Original, MSQ drill]}{\textbf{(MSQ: select ALL correct options.)} Which are stakeholders in IT governance?}
  {\item Senior management
   \item IT staff
   \item System users
   \item External auditors or regulators where applicable}
  {A, B, C, D}{Governance involves organizational leaders, IT teams, users, and relevant external oversight stakeholders.}
\fi
\ifnum\JKSSBLessonID=48
\JKSSBItem{[PYQ-adapted: JKSSB Sub-Inspector 2022, QCRIAS item 391; SRC-17/SRC-22]}{What happens to data held in volatile RAM when power is removed?}
  {\item It is lost
   \item It becomes permanent ROM
   \item It is printed automatically
   \item It is sent to a network router}
  {A}{Volatile RAM loses its contents without power. It is not automatically converted, printed, or transmitted.}
\JKSSBItem{[PYQ-adapted: JKSSB FAA 2022, QCRIAS item 551; SRC-17/SRC-21]}{Which memory is traditionally used to store BIOS firmware?}
  {\item RAM
   \item ROM
   \item Cache
   \item Register}
  {B}{Firmware is stored in non-volatile ROM/flash. RAM, cache, and registers are volatile working storage.}
\JKSSBItem{[PYQ-adapted: JKSSB Website Operator 2026 Q70, provisional key; SRC-04/SRC-07]}{Which method reads magnetic-ink characters on bank cheques?}
  {\item MICR
   \item OCR
   \item OMR
   \item Barcode scanner}
  {A}{MICR reads magnetic-ink characters. OCR recognizes printed text, OMR detects marks, and barcode scanners read barcodes.}
\JKSSBItem{[Original]}{Which shortcut commonly pastes copied content?}
  {\item Ctrl+V
   \item Ctrl+C
   \item Ctrl+X
   \item Ctrl+P}
  {A}{Ctrl+V pastes. Ctrl+C copies, Ctrl+X cuts, and Ctrl+P opens Print.}
\JKSSBItem{[Original, MSQ drill]}{\textbf{(MSQ: select ALL correct options.)} Which pairs are correctly matched?}
  {\item RAM --- volatile memory
   \item MICR --- magnetic-ink character reading
   \item Ctrl+V --- paste
   \item Router --- printer output device}
  {A, B, C}{The first three pairs are correct. A router forwards network traffic; it is not a printer.}
\fi
\ifnum\JKSSBLessonID=49
\JKSSBItem{[Original]}{First-generation electronic computers primarily used:}
  {\item Vacuum tubes
   \item Integrated circuits
   \item Microprocessors
   \item Transistors only}
  {A}{First-generation computers used vacuum tubes. Transistors, integrated circuits, and microprocessors characterize later generations.}
\JKSSBItem{[Original]}{Integrated circuits are most closely associated with which generation?}
  {\item First
   \item Second
   \item Third
   \item Fifth}
  {C}{Third-generation computers used integrated circuits. First used vacuum tubes, second transistors, and fourth microprocessors.}
\JKSSBItem{[Original]}{The microprocessor is most closely associated with which computer generation?}
  {\item First
   \item Second
   \item Third
   \item Fourth}
  {D}{Fourth-generation computers are associated with microprocessors. The other options correspond to earlier technologies.}
\JKSSBItem{[Original]}{Which technology replaced vacuum tubes in second-generation computers?}
  {\item Transistors
   \item Integrated circuits
   \item Microprocessors
   \item Optical fibers}
  {A}{Transistors replaced vacuum tubes in second-generation systems. Integrated circuits came later.}
\JKSSBItem{[Original, MSQ drill]}{\textbf{(MSQ: select ALL correct options.)} Match technology to generation.}
  {\item Vacuum tubes --- first
   \item Transistors --- second
   \item Integrated circuits --- third
   \item Microprocessors --- fourth}
  {A, B, C, D}{All four are the standard technology-generation associations.}
\fi
\ifnum\JKSSBLessonID=50
\JKSSBItem{[Original]}{What is the decimal value of binary 1010?}
  {\item 8
   \item 10
   \item 12
   \item 14}
  {B}{Binary 1010 equals 8+2=10 in decimal. The other options do not match its place-value sum.}
\JKSSBItem{[Original]}{In hexadecimal notation, the digit A represents decimal:}
  {\item 8
   \item 9
   \item 10
   \item 11}
  {C}{Hexadecimal digits A through F represent decimal 10 through 15.}
\JKSSBItem{[Original]}{BCD represents each decimal digit using:}
  {\item Four bits
   \item One byte
   \item Three bits
   \item Sixteen bits}
  {A}{Binary-coded decimal encodes each decimal digit separately in four bits.}
\JKSSBItem{[Original]}{Which logic gate outputs 1 only when its two inputs differ?}
  {\item XOR
   \item AND
   \item OR
   \item NOT}
  {A}{XOR is 1 when inputs differ. AND requires both inputs 1; OR requires at least one; NOT has one input.}
\JKSSBItem{[Original, MSQ drill]}{\textbf{(MSQ: select ALL correct options.)} Which statements are correct?}
  {\item One hexadecimal digit represents four binary bits.
   \item A byte contains eight bits.
   \item BCD encodes each decimal digit in four bits.
   \item Decimal 10 is binary 1110.}
  {A, B, C}{The first three are correct. Binary 1110 equals decimal 14; decimal 10 is binary 1010.}
\fi
\ifnum\JKSSBLessonID=51
\JKSSBItem{[Original]}{Which translator converts an entire high-level program before execution?}
  {\item Compiler
   \item Interpreter
   \item Text editor
   \item Linker only}
  {A}{A compiler translates a program into another form before execution. An interpreter executes through translation as it runs; an editor edits text; a linker combines object modules.}
\JKSSBItem{[Original]}{Which data structure follows the LIFO principle?}
  {\item Stack
   \item Queue
   \item Array only
   \item Graph}
  {A}{A stack is last-in, first-out. A queue is FIFO; arrays and graphs are not defined by LIFO behavior.}
\JKSSBItem{[Original]}{A written step-by-step procedure for solving a problem is an:}
  {\item Algorithm
   \item Operating system
   \item Database record
   \item Device driver}
  {A}{An algorithm is a finite procedure for solving a problem. The other terms describe software/system/data objects.}
\JKSSBItem{[Original]}{Which language is closest to the processor's native instruction representation?}
  {\item Machine language
   \item High-level language
   \item Natural language
   \item Spreadsheet formula language only}
  {A}{Machine language is represented in processor instructions. High-level languages are more abstract; natural language is not processor code.}
\JKSSBItem{[Original, MSQ drill]}{\textbf{(MSQ: select ALL correct options.)} Which pairings are correct?}
  {\item Stack --- LIFO
   \item Queue --- FIFO
   \item Compiler --- translates a program before execution
   \item Debugging --- intentionally adds errors}
  {A, B, C}{The first three are correct. Debugging identifies and fixes errors rather than adding them.}
\fi
\ifnum\JKSSBLessonID=52
\JKSSBItem{[PYQ-adapted: JKSSB Sub-Inspector 2022, QCRIAS item 329; SRC-17/SRC-22]}{Which software category controls computer operations and coordinates hardware?}
  {\item System software
   \item Application software
   \item User data
   \item Presentation theme}
  {A}{System software, including the OS, manages resources and hardware. Application software performs user tasks; data and themes are not that software category.}
\JKSSBItem{[Original]}{In Linux, which command commonly lists files in a directory?}
  {\item ls
   \item cd
   \item pwd
   \item rm}
  {A}{`ls` lists directory contents. `cd` changes directory, `pwd` prints the current path, and `rm` removes files.}
\JKSSBItem{[Original]}{Which Windows utility is intended to remove selected temporary files and free disk space?}
  {\item Disk Cleanup
   \item Paint
   \item Notepad
   \item Character Map}
  {A}{Disk Cleanup removes selected temporary/unneeded files. The other utilities edit/draw text or inspect characters.}
\JKSSBItem{[Original]}{Which type of operating system is designed to respond within defined timing constraints?}
  {\item Real-time operating system
   \item Batch-only database
   \item Spreadsheet system
   \item Web browser}
  {A}{An RTOS is designed for bounded timing behavior. The other options are not operating-system classes.}
\JKSSBItem{[Original, MSQ drill]}{\textbf{(MSQ: select ALL correct options.)} Match the command to its common Linux role.}
  {\item ls --- list directory entries
   \item pwd --- show current working directory
   \item cd --- change directory
   \item rm --- rename a directory only}
  {A, B, C}{The first three are common command roles. `rm` removes files/directories; renaming is commonly done with `mv`.}
\fi
\ifnum\JKSSBLessonID=53
\JKSSBItem{[PYQ-adapted: JKSSB JKP Constable 2024, QCRIAS item 765; SRC-02]}{Which organization develops standards for the World Wide Web?}
  {\item W3C
   \item WHO
   \item OPEC
   \item UNESCO}
  {A}{W3C is the World Wide Web Consortium, which develops Web standards. The other organizations have different mandates.}
\JKSSBItem{[PYQ-adapted: JKSSB Sub-Inspector 2022, QCRIAS item 399; SRC-17/SRC-22]}{Which protocol is commonly used by an email client to retrieve messages from a remote server?}
  {\item SMTP
   \item POP3
   \item SNMP
   \item UDP}
  {B}{POP3 retrieves email. SMTP sends/transfers it; SNMP manages network devices; UDP is a transport protocol.}
\JKSSBItem{[Original]}{Which Internet service transfers files between a client and server?}
  {\item FTP
   \item DNS
   \item ARP
   \item ICMP}
  {A}{FTP is a file-transfer protocol. DNS resolves names, ARP resolves local addresses, and ICMP carries control/diagnostic messages.}
\JKSSBItem{[Original]}{Which legacy protocol provides remote terminal access without encryption by default?}
  {\item Telnet
   \item HTTPS
   \item SFTP
   \item IMAPS}
  {A}{Telnet is a legacy plaintext remote-terminal protocol. HTTPS, SFTP, and IMAPS are secure protocol variants/services.}
\JKSSBItem{[Original, MSQ drill]}{\textbf{(MSQ: select ALL correct options.)} Match protocol/service to its common role.}
  {\item DNS --- name resolution
   \item SMTP --- email sending/transfer
   \item FTP --- file transfer
   \item ICMP --- web page markup}
  {A, B, C}{DNS, SMTP, and FTP match their roles. ICMP is for network control/diagnostics, not HTML markup.}
\fi
\ifnum\JKSSBLessonID=54
\JKSSBItem{[Original]}{Which shortcut commonly opens the Print dialog in Office applications?}
  {\item Ctrl+P
   \item Ctrl+H
   \item Ctrl+K
   \item Ctrl+Z}
  {A}{Ctrl+P opens Print. Ctrl+H opens Find/Replace, Ctrl+K inserts a hyperlink, and Ctrl+Z undoes.}
\JKSSBItem{[Original]}{Which Word shortcut commonly inserts a hyperlink?}
  {\item Ctrl+K
   \item Ctrl+L
   \item Ctrl+M
   \item Ctrl+W}
  {A}{Ctrl+K opens the hyperlink action in common Office versions. The other shortcuts have different or application-dependent uses.}
\JKSSBItem{[Original]}{Which key starts a PowerPoint slide show from the current slide?}
  {\item Shift+F5
   \item F5
   \item F7
   \item Ctrl+P}
  {A}{Shift+F5 starts from the current slide; F5 starts from the beginning. F7 and Ctrl+P are not that action.}
\JKSSBItem{[Original]}{Which key commonly starts spelling and grammar review in Word?}
  {\item F7
   \item F5
   \item F1
   \item F12}
  {A}{F7 invokes spelling/grammar review in common Word versions. F1 is Help, F5 is navigation/go-to, and F12 opens Save As.}
\JKSSBItem{[Original, MSQ drill]}{\textbf{(MSQ: select ALL correct options.)} Match common shortcuts to actions.}
  {\item Ctrl+S --- save
   \item Ctrl+P --- print
   \item Ctrl+Z --- undo
   \item Ctrl+C --- paste}
  {A, B, C}{Ctrl+S saves, Ctrl+P prints, Ctrl+Z undoes, and Ctrl+C copies. Ctrl+V pastes.}
\fi
\ifnum\JKSSBLessonID=55
\JKSSBItem{[PYQ-adapted: JKSSB FAA 2022, QCRIAS item 554; SRC-21]}{Which listed device is a flash-memory device?}
  {\item Hard disk drive
   \item RAM
   \item USB drive
   \item Blu-ray disc}
  {C}{A USB drive uses flash memory. HDDs use magnetic storage, RAM is working memory, and Blu-ray is optical.}
\JKSSBItem{[Original]}{Which extension is commonly associated with a PowerPoint presentation?}
  {\item .pptx
   \item .xlsx
   \item .docx
   \item .accdb}
  {A}{.pptx is PowerPoint; .xlsx is Excel, .docx is Word, and .accdb is Access.}
\JKSSBItem{[Original]}{Which statement correctly distinguishes a file extension from a file format?}
  {\item `.png` is an extension; PNG is the format name.
   \item PNG is an extension and `.png` is the format.
   \item They always indicate the file's actual contents with certainty.
   \item An extension is a storage device.}
  {A}{The dot-suffix in a filename is an extension; PNG names the format. Extensions can be misleading and do not guarantee contents.}
\JKSSBItem{[Original]}{Which file type is typically plain-text tabular data and does not preserve rich workbook formatting?}
  {\item CSV
   \item XLSX
   \item PPTX
   \item DOCX}
  {A}{CSV stores delimited text data. XLSX is a workbook format; PPTX and DOCX are presentation and word-processing formats.}
\JKSSBItem{[Original, MSQ drill]}{\textbf{(MSQ: select ALL correct options.)} Which pairs are correctly matched?}
  {\item .docx --- Word document
   \item .xlsx --- Excel workbook
   \item .pptx --- PowerPoint presentation
   \item .png --- database file}
  {A, B, C}{The first three are correct. PNG is an image format, not a database format.}
\fi
\ifnum\JKSSBLessonID=56
\JKSSBItem{[PYQ-adapted: JKSSB FAA 2022, QCRIAS item 556; SRC-21]}{Which virus type attaches itself to executable files such as .EXE or .COM files?}
  {\item Direct-action file virus
   \item Boot-sector virus
   \item Macro virus
   \item Overwrite virus only}
  {A}{A direct-action file virus infects executable files. Boot-sector viruses target boot records; macro viruses use document macros.}
\JKSSBItem{[PYQ-adapted: JKSSB FAA 2024, QCRIAS item 664; SRC-16/SRC-19]}{Which option is antivirus software?}
  {\item Trojan horse
   \item Worm
   \item Macro virus
   \item Norton}
  {D}{Norton is an antivirus product. A Trojan, worm, and macro virus are malware.}
\JKSSBItem{[Original]}{Which security control filters network traffic according to rules?}
  {\item Firewall
   \item Spreadsheet
   \item Defragmenter
   \item OCR}
  {A}{A firewall filters network traffic. The other choices are an application, storage utility, and text-recognition technology.}
\JKSSBItem{[Original]}{A file that appears harmless but contains a hidden malicious purpose is a:}
  {\item Trojan horse
   \item Patch
   \item Backup
   \item Firewall}
  {A}{A Trojan disguises malicious behavior as benign software. Patches, backups, and firewalls are defensive measures.}
\JKSSBItem{[Original, MSQ drill]}{\textbf{(MSQ: select ALL correct options.)} Which are security controls?}
  {\item Antivirus software
   \item Firewall
   \item Multi-factor authentication
   \item Phishing email}
  {A, B, C}{Antivirus, firewall, and MFA are controls. A phishing email is an attack vector.}
\fi
\ifnum\JKSSBLessonID=57
\JKSSBItem{[PYQ-adapted: JKSSB FAA 2024, QCRIAS item 672; SRC-16/SRC-19]}{The UMANG app enables citizens to access:}
  {\item Virtual-reality games
   \item Streaming movies
   \item Government services and information
   \item Online fitness classes}
  {C}{UMANG is a government-service access platform. The other choices describe entertainment or fitness apps.}
\JKSSBItem{[PYQ-adapted: JKSSB FAA 2024, QCRIAS item 671; SRC-16/SRC-19]}{Smart Governance primarily involves:}
  {\item Making governance processes more complicated
   \item Integrating IT to improve service delivery and decisions
   \item Avoiding technology
   \item Reducing transparency}
  {B}{Smart Governance uses technology to improve efficiency, transparency, and service delivery. The other choices contradict those aims.}
\JKSSBItem{[Original]}{A service that lets a citizen access services from central and state agencies through one mobile platform is an example of:}
  {\item Integrated e-governance delivery
   \item Disk defragmentation
   \item Local file indexing
   \item Spreadsheet recalculation}
  {A}{Integrated e-governance platforms provide access to services across agencies. The other options are local computer operations.}
\JKSSBItem{[Original]}{A Government-to-Citizen service is abbreviated:}
  {\item G2C
   \item G2G
   \item G2B
   \item G2E}
  {A}{G2C means Government to Citizen. G2G is government-to-government, G2B government-to-business, and G2E government-to-employee.}
\JKSSBItem{[Original, MSQ drill]}{\textbf{(MSQ: select ALL correct options.)} Which are common e-governance interaction categories?}
  {\item G2C
   \item G2G
   \item G2B
   \item B2C means Government to Citizen}
  {A, B, C}{G2C, G2G, and G2B are standard categories. B2C means Business to Consumer, not Government to Citizen.}
\fi
\ifnum\JKSSBLessonID=58
\JKSSBItem{[PYQ-adapted: JKSSB Website Operator 2026 Q78, provisional key; SRC-04/SRC-07]}{What character must begin a standard Excel formula?}
  {\item =
   \item \#
   \item @
   \item \&}
  {A}{Excel formulas normally begin with an equals sign. The other characters do not serve as the standard formula prefix.}
\JKSSBItem{[PYQ-adapted: JKSSB Junior Assistant 2026 Q78, final key; SRC-03/SRC-06]}{How many bits are in an IPv6 address?}
  {\item 32
   \item 64
   \item 128
   \item 256}
  {C}{IPv6 addresses are 128 bits. IPv4 uses 32 bits; 64 and 256 are not the address width.}
\JKSSBItem{[Original]}{Which of the following is NOT a storage device?}
  {\item Keyboard
   \item SSD
   \item HDD
   \item USB flash drive}
  {A}{A keyboard is an input device. SSD, HDD, and USB flash drive store data.}
\JKSSBItem{[Original]}{Which action best protects against losing files after accidental deletion or hardware failure?}
  {\item Maintain and test backups
   \item Increase monitor brightness
   \item Rename the desktop folder
   \item Disable all updates}
  {A}{Backups provide recovery copies and should be tested. The other actions do not protect data and disabling updates can increase risk.}
\JKSSBItem{[Original, MSQ drill]}{\textbf{(MSQ: select ALL correct options.)} Select the accurate statements.}
  {\item RAM is volatile.
   \item IPv6 addresses are 128 bits.
   \item Excel formulas commonly begin with `=`.
   \item A section break is the same as a page break.}
  {A, B, C}{The first three statements are correct. A section break creates a new section; a page break alone does not.}
\fi
\ifnum\JKSSBLessonID=59
\JKSSBItem{[PYQ-adapted: SSC CGL 2025 Tier-II Paper-I, 19-Jan-2026 Q6; SRC-18]}{Which register holds the instruction currently being decoded and executed?}
  {\item MAR
   \item PC
   \item IR
   \item Accumulator}
  {C}{The Instruction Register holds the current instruction. MAR stores an address, PC the next instruction address, and the accumulator an intermediate result.}
\JKSSBItem{[PYQ-adapted: JKSSB Junior Assistant 2026 Q65, final key; SRC-03/SRC-06]}{In BCD, each decimal digit is represented using:}
  {\item Four bits
   \item Eight bits
   \item One bit
   \item Sixteen bits}
  {A}{BCD encodes each decimal digit in four bits. It is not the same as converting the whole number to pure binary.}
\JKSSBItem{[PYQ-adapted: JKSSB Junior Assistant 2026 Q65, final key; SRC-03/SRC-06]}{For unsigned widening, which method fills the newly added high-order bits with zeros?}
  {\item Zero extension
   \item Sign extension
   \item Sign inversion
   \item Carry propagation}
  {A}{Zero extension pads an unsigned value with zeros. Sign extension copies the sign bit for signed values.}
\JKSSBItem{[Original]}{In a pipelined processor, a data hazard can occur when:}
  {\item An instruction depends on a result not yet produced by an earlier instruction
   \item A printer runs out of paper
   \item A file extension is hidden
   \item A keyboard key is repeated}
  {A}{A data hazard is a dependency conflict between overlapping instructions. The other choices are peripheral or file-management issues.}
\JKSSBItem{[Original, MSQ drill]}{\textbf{(MSQ: select ALL correct options.)} Which are standard CPU architecture concepts?}
  {\item Program Counter stores the next instruction address.
   \item Instruction Register holds the current instruction.
   \item A carry-out and signed overflow are always identical.
   \item BCD uses four bits per decimal digit.}
  {A, B, D}{PC, IR, and BCD statements are correct. Carry-out is not generally equivalent to signed overflow.}
\fi
\ifnum\JKSSBLessonID=60
\JKSSBItem{[PYQ-adapted: JKSSB Junior Assistant 2026 Q70, final key; SRC-03/SRC-06]}{Which deadlock recovery method takes a resource away from a process to resolve a resource cycle?}
  {\item Resource preemption
   \item Increasing display resolution
   \item Formatting the keyboard
   \item Renaming a user account}
  {A}{Resource preemption reclaims a resource to break deadlock. The other actions do not change process-resource allocation.}
\JKSSBItem{[Original]}{In Linux, which command prints the current working directory?}
  {\item pwd
   \item ls
   \item cd
   \item rm}
  {A}{`pwd` prints the current directory. `ls` lists entries, `cd` changes directory, and `rm` removes files.}
\JKSSBItem{[Original]}{Which command is commonly used to list directory contents in Linux?}
  {\item ls
   \item pwd
   \item mkdir
   \item rmdir}
  {A}{`ls` lists directory entries. `pwd` prints the path; `mkdir` and `rmdir` create/remove directories.}
\JKSSBItem{[Original]}{Which scheduling policy gives each ready process a fixed time slice in turn?}
  {\item Round-robin
   \item First-come, first-served only
   \item Shortest job first only
   \item Non-preemptive batch sorting}
  {A}{Round-robin cycles through ready processes using a time quantum. The other options do not specify that rotating time-slice policy.}
\JKSSBItem{[Original, MSQ drill]}{\textbf{(MSQ: select ALL correct options.)} Which are common operating-system concepts?}
  {\item A process is a program in execution.
   \item Deadlock involves processes waiting on resources.
   \item A file system organizes stored files.
   \item A device driver is a physical keyboard.}
  {A, B, C}{Processes, deadlocks, and file systems are OS concepts. A driver is software that helps control a device.}
\fi
\ifnum\JKSSBLessonID=61
\JKSSBItem{[PYQ-adapted: SSC CGL 2025 Tier-II Paper-I, 19-Jan-2026 Q14; SRC-18]}{Which setting can be changed for one Word section without changing the previous section?}
  {\item Page orientation and margins
   \item AutoCorrect replacement list
   \item Normal style definition
   \item File Properties metadata}
  {A}{Page setup is section-specific. AutoCorrect is application-wide; styles and Properties are not unique page-layout settings for one section.}
\JKSSBItem{[PYQ-adapted: SSC CGL 2025 Tier-II Paper-I, 19-Jan-2026 Q4; SRC-18]}{Which Excel function counts cells that meet one criterion?}
  {\item COUNT
   \item COUNTBLANK
   \item COUNTIF
   \item SUMIF}
  {C}{COUNTIF counts matching cells. COUNT counts numeric cells without a criterion; COUNTBLANK counts empties; SUMIF sums matching values.}
\JKSSBItem{[PYQ-adapted: JKSSB Junior Assistant 2026 Q76, final key; SRC-03/SRC-06]}{Which email field hides listed recipients from other recipients?}
  {\item To
   \item Cc
   \item Bcc
   \item Subject}
  {C}{Bcc conceals those recipients from other recipients. To and Cc recipients are visible; Subject identifies the topic.}
\JKSSBItem{[Original]}{Which feature in Word records edits for later review?}
  {\item Track Changes
   \item Mail Merge
   \item AutoSum
   \item Slide Sorter}
  {A}{Track Changes records revisions. Mail Merge personalizes documents; AutoSum is an Excel feature; Slide Sorter is a PowerPoint view.}
\JKSSBItem{[Original, MSQ drill]}{\textbf{(MSQ: select ALL correct options.)} Which distinctions are accurate?}
  {\item Section Break can enable different page setup for a new section.
   \item Bcc conceals recipients from other recipients.
   \item COUNTIF counts cells meeting criteria.
   \item Digital signatures are the same as message encryption.}
  {A, B, C}{Section setup, Bcc, and COUNTIF are correctly described. Digital signatures support authenticity/integrity; encryption supports confidentiality.}
\fi
\ifnum\JKSSBLessonID=62
\JKSSBItem{[PYQ-adapted: JKSSB Junior Assistant 2026 Q78, final key; SRC-03/SRC-06]}{How many bits are in an IPv6 address?}
  {\item 32
   \item 64
   \item 128
   \item 256}
  {C}{IPv6 addresses are 128 bits. IPv4 addresses are 32 bits.}
\JKSSBItem{[PYQ-adapted: SSC CGL 2025 Tier-II Paper-I, 19-Jan-2026 Q3; SRC-18]}{Which OSI layer routes packets across networks using logical addressing?}
  {\item Data Link
   \item Session
   \item Network
   \item Transport}
  {C}{The Network layer provides logical addressing and routing. Data Link handles local frames; Session manages dialogues; Transport provides end-to-end services.}
\JKSSBItem{[PYQ-adapted: JKSSB Junior Assistant 2026 Q79, final key; SRC-03/SRC-06]}{HTTP/3 uses which transport protocol through QUIC?}
  {\item UDP
   \item TCP
   \item ICMP
   \item ARP}
  {A}{HTTP/3 uses QUIC over UDP. TCP is used by earlier HTTP versions; ICMP and ARP have other network roles.}
\JKSSBItem{[Original]}{Which cloud service model gives a customer virtual machines, storage, and networking while the customer manages the operating system?}
  {\item IaaS
   \item SaaS
   \item PaaS
   \item BIOS}
  {A}{IaaS provides infrastructure resources; the customer manages the guest OS and applications. SaaS delivers an application; PaaS provides a managed platform.}
\JKSSBItem{[Original, MSQ drill]}{\textbf{(MSQ: select ALL correct options.)} Which statements are correct?}
  {\item IPv6 uses 128-bit addresses.
   \item Routers forward packets between networks.
   \item UDP is connectionless.
   \item HTTP/3 uses TCP through QUIC.}
  {A, B, C}{The first three are correct. QUIC runs over UDP, not TCP.}
\fi
\ifnum\JKSSBLessonID=63
\JKSSBItem{[PYQ-adapted: JKSSB Sub-Inspector 2022, QCRIAS item 359; SRC-17/SRC-22]}{Which markup language is commonly used to add dynamic effects to a web page in older exam terminology?}
  {\item DHTML
   \item XML only
   \item C language
   \item Assembly language}
  {A}{DHTML describes combining HTML, scripting, and style changes for dynamic page behavior. The other choices do not name that presentation concept.}
\JKSSBItem{[Original]}{Which language primarily structures the content of a web page?}
  {\item HTML
   \item SMTP
   \item SQL only
   \item BIOS}
  {A}{HTML marks up web content. SMTP transfers email, SQL queries databases, and BIOS is firmware.}
\JKSSBItem{[Original]}{Which technology is primarily used to add client-side interactivity to a web page?}
  {\item JavaScript
   \item JPEG
   \item FTP
   \item ROM}
  {A}{JavaScript is commonly used for client-side behavior. JPEG is an image format, FTP transfers files, and ROM is memory.}
\JKSSBItem{[Original]}{XML is primarily designed to:}
  {\item Mark up and structure data
   \item Route IP packets
   \item Replace all programming languages
   \item Print magnetic-ink characters}
  {A}{XML structures and describes data. Packet routing, replacing all languages, and MICR printing are unrelated.}
\JKSSBItem{[Original, MSQ drill]}{\textbf{(MSQ: select ALL correct options.)} Which are web-related technologies/languages?}
  {\item HTML
   \item JavaScript
   \item XML
   \item SMTP}
  {A, B, C}{HTML, JavaScript, and XML are web technologies/languages. SMTP is an email-transfer protocol.}
\fi
\ifnum\JKSSBLessonID=64
\JKSSBItem{[PYQ-adapted: SSC CGL 2025 Tier-II Paper-I, 19-Jan-2026 Q5; SRC-18]}{Which option provides message-level email encryption to protect content from unauthorized access?}
  {\item S/MIME
   \item IMAP
   \item SSL alone
   \item SMTP}
  {A}{S/MIME can encrypt the message itself. IMAP retrieves mail, SMTP transfers it, and TLS/SSL on one connection hop is transport protection rather than end-to-end message encryption.}
\JKSSBItem{[PYQ-adapted: SSC CGL 2025 Tier-II Paper-I, 19-Jan-2026 Q11; SRC-18]}{What is a common purpose of a Trojan horse?}
  {\item Disguise malicious activity as a benign program
   \item Self-replicate as a worm
   \item Filter network packets as a firewall
   \item Repair operating-system vulnerabilities}
  {A}{A Trojan hides a malicious purpose inside something that appears benign. Self-replication describes a worm; the other options describe defensive controls.}
\JKSSBItem{[Original]}{Which property does encryption primarily provide when only authorized parties can read the content?}
  {\item Confidentiality
   \item Availability
   \item Non-repudiation only
   \item Compression}
  {A}{Encryption protects confidentiality. Availability concerns access; signatures can support non-repudiation/integrity; compression reduces size.}
\JKSSBItem{[Original]}{A digital signature primarily helps verify:}
  {\item Authenticity and integrity
   \item Secrecy of the message content by itself
   \item Internet connection speed
   \item File size reduction}
  {A}{A digital signature supports sender authenticity and message integrity. Encryption, not signing alone, provides confidentiality.}
\JKSSBItem{[Original, MSQ drill]}{\textbf{(MSQ: select ALL correct options.)} Match security concept to purpose.}
  {\item Encryption --- confidentiality
   \item Digital signature --- authenticity/integrity
   \item Firewall --- traffic control
   \item Antivirus --- guaranteed prevention of every attack}
  {A, B, C}{The first three are standard roles. Antivirus reduces risk but cannot guarantee prevention of every attack.}
\fi
\ifnum\JKSSBLessonID=65
\JKSSBItem{[PYQ-adapted: JKSSB Website Operator 2026 Q69, provisional key; SRC-04/SRC-07]}{Which device is used to produce large engineering drawings?}
  {\item Plotter
   \item Microphone
   \item Scanner
   \item Trackball}
  {A}{A plotter produces large-format output. The other devices provide input.}
\JKSSBItem{[PYQ-adapted: JKSSB Junior Assistant 2026 Q68, final key; SRC-03/SRC-06]}{Which device is an output device?}
  {\item Plotter
   \item Scanner
   \item Light pen
   \item Digitizer tablet}
  {A}{A plotter outputs drawings. A scanner, light pen, and digitizer tablet provide input.}
\JKSSBItem{[Original]}{Sampling in digital audio converts a continuous signal into:}
  {\item Measurements taken at discrete time intervals
   \item A printed paper page
   \item A file-system directory
   \item A network address}
  {A}{Sampling records signal values at discrete times. The other choices are unrelated.}
\JKSSBItem{[Original]}{Which file format commonly stores compressed still images?}
  {\item JPEG
   \item DOCX
   \item CSV
   \item EXE}
  {A}{JPEG is an image format commonly using lossy compression. DOCX is a document package, CSV is delimited text, and EXE is executable.}
\JKSSBItem{[Original, MSQ drill]}{\textbf{(MSQ: select ALL correct options.)} Which are multimedia output devices?}
  {\item Monitor
   \item Speaker
   \item Projector
   \item Scanner}
  {A, B, C}{Monitor, speaker, and projector output visual/audio content. A scanner captures input.}
\fi
\ifnum\JKSSBLessonID=66
\JKSSBItem{[Original]}{A question copied from a past paper should be labeled as a PYQ only when:}
  {\item Its exam identity and source are traceable
   \item It resembles a familiar exam style
   \item It appears in an uncredited mock test
   \item Its answer seems obvious}
  {A}{A PYQ claim requires a traceable paper identity/source. Similar style or plausibility does not establish provenance.}
\JKSSBItem{[Original]}{If a secondary transcription's answer conflicts with the official final key, which should govern the answer?}
  {\item Official final key
   \item Secondary transcription
   \item A generated explanation
   \item Whichever option appears first}
  {A}{The official final answer key is authoritative over a secondary transcription or generated explanation.}
\JKSSBItem{[Original]}{An MSQ written for practice based on a single-answer exam topic should be labeled:}
  {\item Original MSQ drill
   \item Verbatim PYQ
   \item Official multi-select question
   \item Final answer key}
  {A}{A newly authored multi-select item is an original MSQ drill. It must not imply the source exam used that format.}
\JKSSBItem{[Original]}{Which information belongs in a reliable PYQ provenance note?}
  {\item Exam/post, year, question identity, and source/key status when known
   \item Only the topic name
   \item The writer's guess about the year
   \item A generated question's difficulty label only}
  {A}{Useful provenance identifies the paper/item and answer-key status. Topic or difficulty alone cannot establish authenticity.}
\JKSSBItem{[Original, MSQ drill]}{\textbf{(MSQ: select ALL correct options.)} Which curation practices protect PYQ integrity?}
  {\item Keep generated questions labeled as original.
   \item Preserve the exact exam/date/question source for a PYQ.
   \item Check disputed answers against the official key.
   \item Label every exam-style question as an authentic PYQ.}
  {A, B, C}{The first three preserve provenance. Exam-style similarity does not make an item an authentic PYQ.}
\fi
\end{enumerate}

\end{document}
